

\setcounter{chapter}{6}

\chapter{环的引论}

\section{基本定义与例子}

\begin{definition}
（1）\textbf{环} \(R\) 是一个集合连同两个二元运算 \(+\) 和 \(\times\)（称为加法和乘法），满足下列公理：

（i）\((R, +)\) 是阿贝尔群；

（ii）\(\times\) 是结合的：对所有 \(a, b, c \in R\) 有 \((a \times b) \times c = a \times (b \times c)\)；

（iii）分配律对 \(R\) 中所有 \(a, b, c \in R\) 成立：
\[
(a+b)\times c = (a\times c) + (b\times c), \qquad a\times(b+c) = (a\times b) + (a\times c).
\]

（2）若乘法交换，则称环 \(R\) 是\textbf{交换的}。

（3）若存在 \(1 \in R\) 使得对所有 \(a \in R\) 有 \(1 \times a = a \times 1 = a\)，则称环 \(R\) 有\textbf{单位元}（或含有一个 1）。
\end{definition}

对 \(a, b \in R\)，我们通常简单地写 \(ab\) 而不是 \(a \times b\). \(R\) 的加法单位元总记作 \(0\)，环元素 \(a\) 的加法逆元记作 \(-a\).

\(R\) 在加法下是群这一条件相当自然，但要求这个群是阿贝尔群可能显得不自然。这样做的一个动机是：若环 \(R\) 含有 1，则加法下的交换性由分配律迫使。为看出这一点，用分配律（但不假设加法交换）以两种方式计算乘积 \((1+1)(a+b)\). 得到
\[
(1+1)(a+b) = 1(a+b) + 1(a+b) = 1a+1b+1a+1b = a+b+a+b
\]
和
\[
(1+1)(a+b) = (1+1)a + (1+1)b = 1a+1a+1b+1b = a+a+b+b.
\]
由于 \(R\) 在加法下是群，这蕴含 \(b+a = a+b\)，即 \(R\) 在加法下必交换。

域是环最重要的例子之一。注意下面给出的定义只是 1.4 节定义的另一种表述。

\begin{definition}
带单位元 \(1\)（\(1 \neq 0\)）的环 \(R\) 称为\textbf{除环}（或\textbf{斜域}），若每个非零元素 \(a \in R\) 都有乘法逆元，即存在 \(b \in R\) 使 \(ab = ba = 1\). 交换的除环称为\textbf{域}。
\end{definition}

下面给出更多环的例子。

\begin{example}
（1）最简单的环的例子是平凡环：取 \(R\) 为任意交换群（群运算记作 \(+\)），并定义 \(R\) 上的乘法 \(\times\) 为对所有 \(a, b \in R\) 有 \(a \times b = 0\). 容易看出这个乘法定义一个交换环。特别地，若 \(R = \{0\}\) 是平凡群，所得环 \(R\) 称为\textbf{零环}，记作 \(R = 0\). 除零环外，平凡环不含单位元（\(R = 0\) 是唯一 \(1 = 0\) 的环；我们将常常通过施加条件 \(1 \neq 0\) 来排除这个环）。虽然平凡环有两个二元运算，乘法并不给加法群增加新结构，环论不提供任何不能由（阿贝尔）群论得到的信息。

（2）整数环 \(\mathbb{Z}\) 在通常的加法与乘法下是带单位元（整数 1）的交换环。环公理（如同加法群公理）由自然数系统的基本公理推出。注意在乘法下 \(\mathbb{Z} - \{0\}\) 不是群（事实上这个环中元素有乘法逆元的极少）。我们不久将回到这些逆元的问题。

（3）类似地，有理数 \(\mathbb{Q}\)、实数 \(\mathbb{R}\) 和复数 \(\mathbb{C}\) 是带单位元的交换环（事实上它们是域）。它们的环公理最终由 \(\mathbb{Z}\) 的环公理推出。我们将在从 \(\mathbb{Z}\) 构造 \(\mathbb{Q}\)（7.5 节）和从 \(\mathbb{R}\) 构造 \(\mathbb{C}\)（13.1 节例 1）时验证这一点；这两种构造都是更一般过程的特殊情形。从 \(\mathbb{Q}\) 构造 \(\mathbb{R}\)（以及随后的环公理验证）在基础分析教材中进行。

（4）商群 \(\mathbb{Z}/n\mathbb{Z}\) 在剩余类的加法与乘法下是带单位元（元素 1）的交换环（常称为“模算术”）。我们已经看到加法阿贝尔群公理由商群理论的一般原理推出（事实上这是商群的原型）。我们不久将证明其余的环公理（特别是剩余类乘法的良定义性）类似地由商环的一般理论推出。

以上所有例子中环都是交换的。历史上最早的非交换环之一由 William Rowan Hamilton 爵士（1805--1865）于 1843 年发现。这个环是除环，对后来数学的发展有极大影响，至今仍在数学和物理的某些领域中起重要作用。

（5）（实 Hamilton 四元数）设 \(\mathbb{H}\) 是形如 \(a + bi + cj + dk\)（\(a, b, c, d \in \mathbb{R}\) 是实数）的元素组成的集合（粗略地说，“系数为实数的 \(1, i, j, k\) 的多项式”），其中加法由
\[
(a + bi + cj + dk) + (a' + b'i + c'j + d'k) = (a+a') + (b+b')i + (c+c')j + (d+d')k
\]
“按分量”定义，乘法通过展开 \((a+bi+cj+dk)(a'+b'i+c'j+d'k)\) 并利用分配律（注意各项的顺序）以及关系
\[
i^2 = j^2 = k^2 = -1, \quad ij = k = -ji,\quad jk = i = -kj,\quad ki = j = -ik
\]
（其中实数系数与 \(i, j, k\) 交换）化简得到。例如，
\[
(1+i+2j)(j+k) = 1(j+k) + i(j+k) + 2j(j+k) = j+k+ij+ik+2j^2+2jk
\]
\[
= j+k+k+(-j)+2(-1)+2(i) = -2+2i+2k.
\]
\(\mathbb{H}\) 是环这一事实可以通过直接但冗长的公理检验来证明（乘法结合律尤其麻烦）。Hamilton 四元数是带单位元（\(1 = 1+0i+0j+0k\)）的非交换环。类似地，取 \(a, b, c, d\) 为有理数可以定义有理 Hamilton 四元数环。实和有理 Hamilton 四元数都是除环，其中非零元素的逆元由
\[
(a+bi+cj+dk)^{-1} = \frac{a - bi - cj - dk}{a^2+b^2+c^2+d^2}
\]
给出。

（6）一类重要的环由考虑函数环得到。设 \(X\) 是任意非空集合，\(A\) 是任意环。所有（集合）函数 \(f : X \to A\) 的集合 \(R\) 在通常的逐点加法与乘法 \((f+g)(x) = f(x)+g(x)\)、\((fg)(x) = f(x)g(x)\) 下是环。\(R\) 的每条环公理都由 \(A\) 的相应公理直接推出。环 \(R\) 交换当且仅当 \(A\) 交换，\(R\) 有 1 当且仅当 \(A\) 有 1（此时 \(R\) 的 1 必是 \(X\) 上的常函数 1）。

若 \(X\) 和 \(A\) 有更多结构，我们可以构造保持这些结构的其他函数环。例如，若 \(A\) 是实数环 \(\mathbb{R}\)，\(X\) 是 \(\mathbb{R}\) 中的闭区间 \([0,1]\)，我们可以构造所有从 \([0,1]\) 到 \(\mathbb{R}\) 的连续函数构成的环（这里需要基本极限定理保证连续函数的和与积连续）——这是带 1 的交换环。

（7）一个没有单位元的环的例子是偶数环 \(2\mathbb{Z}\)，在通常的整数加法与乘法下（偶数的和与积仍是偶数）。

分析中自然出现的另一个例子如下。函数 \(f : \mathbb{R} \to \mathbb{R}\) 称为\textbf{具有紧支集}，若存在实数 \(a, b\)（依赖于 \(f\)）使对所有 \(x \notin [a,b]\) 有 \(f(x) = 0\)（即 \(f\) 在某个有界区间外为零）。所有具有紧支集的函数 \(f : \mathbb{R} \to \mathbb{R}\) 的集合是无单位元的交换环（因为单位元不可能具有紧支集）。类似地，所有具有紧支集的连续函数 \(f : \mathbb{R} \to \mathbb{R}\) 的集合也是无单位元的交换环。
\end{example}

下一节我们给出三种从给定环构造“较大”环的重要方式（类比于上面的例 6），从而大大扩充我们的例子列表。在这样做之前，我们提到任意环的一些基本性质。

环 \(\mathbb{Z}\) 是一个值得记住的好例子，尽管这个环比一般环有更多的代数结构（例如它交换且有单位元）。然而，它的基本算术对一般环也成立，如下面的结果所示。

\begin{proposition}\label{pro:7.1}
设 \(R\) 是环。则

\begin{enumerate}[label=(\arabic*)]
\item 对所有 \(a \in R\) 有 \(0a = a0 = 0\).
\item 对所有 \(a, b \in R\) 有 \((-a)b = a(-b) = -(ab)\)（回忆 \(-a\) 是 \(a\) 的加法逆元）。
\item 对所有 \(a, b \in R\) 有 \((-a)(-b) = ab\).
\item 若 \(R\) 有单位元 1，则单位元唯一，且 \(-a = (-1)a\).
\end{enumerate}
\end{proposition}

\begin{proof}
这些都由分配律和加法群 \(R\) 中的消去律推出。例如，（1）由 \(0a = (0+0)a = 0a + 0a\) 推出。（2）中的等式 \((-a)b = -(ab)\) 由 \(ab + (-a)b = (a + (-a))b = 0b = 0\) 推出。其余部分类似地推出，留作读者练习。
\end{proof}

这个命题表明，由于分配律，环的加法与乘法结构彼此配合良好，正如熟悉的整数例子一样。

然而与整数不同，一般环可能有许多元素有乘法逆元，也可能有非零元素 \(a, b\) 使 \(ab = 0\). 元素的这两个与环的乘法结构有关的性质有专门的名称。

\begin{definition}
设 \(R\) 是环。

（1）\(R\) 的非零元素 \(a\) 称为\textbf{零因子}，若存在 \(R\) 中的非零元素 \(b\) 使 \(ab = 0\) 或 \(ba = 0\).

（2）假设 \(R\) 有单位元 \(1 \neq 0\). \(R\) 的元素 \(u\) 称为 \(R\) 中的\textbf{单位}（可逆元），若存在 \(R\) 中的 \(v\) 使 \(uv = vu = 1\). \(R\) 中单位的集合记作 \(R^\times\).
\end{definition}

容易看出环 \(R\) 中的单位在乘法下构成群，故 \(R^\times\) 称为 \(R\) 的\textbf{单位群}。用这个术语，域是带单位元 \(1 \neq 0\)、其中每个非零元素都是单位的交换环 \(F\)，即 \(F^\times = F - \{0\}\).

注意零因子永远不可能是单位。例如假设 \(a\) 是 \(R\) 中的单位，且对某个非零 \(b \in R\) 有 \(ab = 0\). 则对某个 \(v \in R\) 有 \(va = 1\)，故 \(b = 1b = (va)b = v(ab) = v0 = 0\)，矛盾。类似地，若对某个非零 \(b\) 有 \(ba = 0\)，则 \(a\) 不可能是单位。这特别表明域不含零因子。

\begin{example}
（1）整数环 \(\mathbb{Z}\) 没有零因子，其唯一的单位是 \(\pm 1\)，即 \(\mathbb{Z}^\times = \{\pm 1\}\). 注意每个非零整数在更大的环 \(\mathbb{Q}\) 中都有逆元，故“是单位”这一性质依赖于元素被看作在哪个环中。

（2）设 \(n \geq 2\) 是整数。在环 \(\mathbb{Z}/n\mathbb{Z}\) 中，与 \(n\) 互素的元素 \(u\) 是单位（我们将在下一章证明这一点）。因此我们使用记号 \((\mathbb{Z}/n\mathbb{Z})^\times\) 与任意环中单位群的定义一致。

另一方面，若 \(a\) 是非零整数且 \(a\) 与 \(n\) 不互素，我们证明 \(\bar{a}\) 是 \(\mathbb{Z}/n\mathbb{Z}\) 中的零因子。为此设 \(d\) 是 \(a\) 与 \(n\) 的最大公因子，\(b = \frac{n}{d}\). 由假设 \(d > 1\)，故 \(0 < b < n\)，即 \(\bar{b} \neq 0\). 但由构造 \(n\) 整除 \(ab\)，即 \(\bar{a}\bar{b} = 0\)（在 \(\mathbb{Z}/n\mathbb{Z}\) 中）。这表明 \(\mathbb{Z}/n\mathbb{Z}\) 的每个非零元素要么是单位要么是零因子。此外，每个非零元素都是单位当且仅当 \(0 < a < n\) 范围内的每个整数 \(a\) 都与 \(n\) 互素。这当且仅当 \(n\) 是素数，即 \(\mathbb{Z}/n\mathbb{Z}\) 是域当且仅当 \(n\) 是素数。

（3）若 \(R\) 是所有从闭区间 \([0,1]\) 到 \(\mathbb{R}\) 的函数构成的环，则 \(R\) 的单位是在任何点都不为零的函数（对这样的 \(f\)，其逆元是函数 \(1/f\)）。若 \(f\) 不是单位且不为零，则 \(f\) 是零因子，因为若定义
\[
g(x) = \begin{cases} 0, & \text{若 } f(x) \neq 0, \\ 1, & \text{若 } f(x) = 0, \end{cases}
\]
则 \(g\) 不是零函数但对所有 \(x\) 有 \(f(x)g(x) = 0\).

（4）若 \(R\) 是所有从闭区间 \([0,1]\) 到 \(\mathbb{R}\) 的连续函数构成的环，则 \(R\) 的单位仍是在任何点都不为零的函数，但现在有既不是单位也不是零因子的函数。例如 \(f(x) = x - \frac{1}{2}\) 只有一个零点（在 \(x = \frac{1}{2}\) 处），故 \(f\) 不是单位。另一方面，若 \(gf = 0\)，则 \(g\) 必须对所有 \(x \neq \frac{1}{2}\) 为零，而具有这一性质的唯一连续函数是零函数。因此 \(f\) 既不是单位也不是零因子。类似地，在 \([0,1]\) 上只有有限（或可数）个零点的函数都不是零因子。这个环也含许多零因子。例如设
\[
f(x) = \begin{cases} 0, & 0 \leq x \leq \frac{1}{2}, \\ x - \frac{1}{2}, & \frac{1}{2} \leq x \leq 1, \end{cases}
\]
并令 \(g(x) = f(1-x)\). 则 \(f\) 和 \(g\) 是非零连续函数，其乘积是零函数。

（5）设 \(D\) 是不是 \(\mathbb{Q}\) 中完全平方的有理数，定义
\[
\mathbb{Q}(\sqrt{D}) = \{a + b\sqrt{D} \mid a, b \in \mathbb{Q}\}
\]
为 \(\mathbb{C}\) 的子集。这个集合显然对减法封闭，且恒等式 \((a+b\sqrt{D})(c+d\sqrt{D}) = (ac+bdD) + (ad+bc)\sqrt{D}\) 表明它对乘法也封闭。因此 \(\mathbb{Q}(\sqrt{D})\) 是 \(\mathbb{C}\) 的子环（若 \(D > 0\) 甚至是 \(\mathbb{R}\) 的子环），特别地是带单位元的交换环。容易证明 \(D\) 不是平方这一假设蕴含 \(\mathbb{Q}(\sqrt{D})\) 的每个元素都可以唯一地写成 \(a + b\sqrt{D}\) 的形式。这个假设还蕴含若 \(a, b\) 不同时为零则 \(a^2 - Db^2 \neq 0\)，且由于 \((a+b\sqrt{D})(a-b\sqrt{D}) = a^2 - Db^2\)，可知若 \(a+b\sqrt{D} \neq 0\)（即 \(a\) 或 \(b\) 非零），则
\[
\frac{a - b\sqrt{D}}{a^2 - Db^2}
\]
是 \(a+b\sqrt{D}\) 在 \(\mathbb{Q}(\sqrt{D})\) 中的逆元。这表明这个交换环中每个非零元素都是单位，即 \(\mathbb{Q}(\sqrt{D})\) 是域（称为\textbf{二次域}，参见 13.2 节）。

有理数 \(D\) 可以写成 \(D = f^2 D'\)，其中 \(f\) 是有理数，\(D'\) 是唯一的不被任何大于 1 的整数的平方整除的整数，即 \(D'\) 是 \(-1\) 或 \(\pm 1\) 乘以 \(\mathbb{Z}\) 中互不相同素数的乘积（例如 \(8/5 = (2/5)^2 \cdot 10\)）。称 \(D'\) 为 \(D\) 的\textbf{无平方因子部分}。则 \(\sqrt{D} = f\sqrt{D'}\)，故 \(\mathbb{Q}(\sqrt{D}) = \mathbb{Q}(\sqrt{D'})\). 因此在二次域 \(\mathbb{Q}(\sqrt{D})\) 的定义中假设 \(D\) 是无平方因子的整数（即 \(f = 1\)）不会失去一般性。
\end{example}

具有与整数 \(\mathbb{Z}\) 某些相同特征的环被赋予一个名称：

\begin{definition}
带单位元 \(1 \neq 0\) 的交换环称为\textbf{整环}，若它没有零因子。
\end{definition}

整环中没有零因子给这些环一个消去性质：

\begin{proposition}\label{pro:7.2}
假设 \(a, b, c\) 是任意环的元素，且 \(a\) 不是零因子。若 \(ab = ac\)，则 \(a = 0\) 或 \(b = c\)（即若 \(a \neq 0\) 我们可以消去 \(a\)）。特别地，若 \(a, b, c\) 是整环中的任意元素且 \(ab = ac\)，则 \(a = 0\) 或 \(b = c\).
\end{proposition}

\begin{proof}
若 \(ab = ac\) 则 \(a(b-c) = 0\)，故 \(a = 0\) 或 \(b - c = 0\). 第二个论断由第一个和整环的定义推出。
\end{proof}

\begin{corollary}\label{cor:7.3}
任何有限整环都是域。
\end{corollary}

\begin{proof}
设 \(R\) 是有限整环，\(a\) 是 \(R\) 的非零元素。由消去律，映射 \(x \mapsto ax\) 是单射。由于 \(R\) 有限，这个映射也是满射。特别地存在 \(b \in R\) 使 \(ab = 1\)，即 \(a\) 是 \(R\) 中的单位。由于 \(a\) 是任意非零元素，\(R\) 是域。
\end{proof}

Wedderburn 的一个著名结果是：有限除环必然交换，即是域。这个定理的证明在 13.6 节末的习题中给出概要。

在 7.5 节我们将更详细地研究零因子与单位的关系。我们将看到交换环中每个不是零因子的非零元素都在某个更大的环中有乘法逆元。这给出了对命题 \ref{pro:7.2} 中消去律的另一种看法。

定义了环的概念之后，自然有子环的概念。

\begin{definition}
环 \(R\) 的\textbf{子环}是 \(R\) 的、对乘法封闭的子群。
\end{definition}

换句话说，环 \(R\) 的子集 \(S\) 是子环，若 \(R\) 中限制到 \(S\) 的加法与乘法给出 \(S\) 的环结构。要证明环 \(R\) 的子集是子环，只需验证它非空且对减法与乘法封闭。

\begin{example}
上面许多例子也是子环。

（1）\(\mathbb{Z}\) 是 \(\mathbb{Q}\) 的子环，\(\mathbb{Q}\) 是 \(\mathbb{R}\) 的子环。“是……的子环”这一性质显然是传递的。

（2）\(2\mathbb{Z}\) 是 \(\mathbb{Z}\) 的子环，对任意整数 \(n\) 的 \(n\mathbb{Z}\) 也是。环 \(\mathbb{Z}/n\mathbb{Z}\) 对任意 \(n \geq 2\) 都不是 \(\mathbb{Z}\) 的子环（也不是子群）。

（3）所有从 \(\mathbb{R}\) 到 \(\mathbb{R}\) 的连续函数构成的环是所有从 \(\mathbb{R}\) 到 \(\mathbb{R}\) 的函数环的子环。所有从 \(\mathbb{R}\) 到 \(\mathbb{R}\) 的可微函数构成的环是这两者的子环。

（4）\(S = \mathbb{Z} + \mathbb{Z}i + \mathbb{Z}j + \mathbb{Z}k\)，即整四元数，构成实或有理四元数的子环——容易验证两个这样的四元数相乘得到另一个系数为整数的四元数。这个环（不是除环）可以用来给出数论中许多结果的证明。

（5）若 \(R\) 是域 \(F\) 的、包含 \(F\) 的单位元的子环，则 \(R\) 是整环。其逆也成立，即任何整环都含于某个域中（参见 7.5 节）。
\end{example}

\begin{example}[二次整数环]
设 \(D\) 是无平方因子的整数。由加法与乘法立即看出集合 \(\mathbb{Z}[\sqrt{D}] = \{a + b\sqrt{D} \mid a, b \in \mathbb{Z}\}\) 构成前面定义的二次域 \(\mathbb{Q}(\sqrt{D})\) 的子环。若 \(D \equiv 1 \pmod 4\)，则稍大的子集
\[
\mathbb{Z}\left[\frac{1+\sqrt{D}}{2}\right] = \left\{a + b\frac{1+\sqrt{D}}{2} \ \middle|\ a, b \in \mathbb{Z}\right\}
\]
也是子环：对加法封闭是显然的，而
\[
\left(a + b\frac{1+\sqrt{D}}{2}\right)\left(c + d\frac{1+\sqrt{D}}{2}\right) = \left(ac + bd\frac{1-D}{4}\right) + \left(ad + bc + bd\right)\frac{1+\sqrt{D}}{2}
\]
连同关于 \(D\) 的同余条件表明对乘法封闭。

定义
\[
\mathcal{O} = \mathcal{O}_{\mathbb{Q}(\sqrt{D})} = \mathbb{Z}[\omega] = \{a + b\omega \mid a, b \in \mathbb{Z}\},
\]
其中
\[
\omega = \begin{cases} \sqrt{D}, & \text{若 } D \equiv 2, 3 \pmod 4, \\ \dfrac{1+\sqrt{D}}{2}, & \text{若 } D \equiv 1 \pmod 4, \end{cases}
\]
称为二次域 \(\mathbb{Q}(\sqrt{D})\) 中的\textbf{整数环}。这一术语来自如下事实：域 \(\mathbb{Q}(\sqrt{D})\) 的子环 \(\mathcal{O}\) 的元素具有许多与有理数域 \(\mathbb{Q}\) 的子环 \(\mathbb{Z}\) 类似的性质（且它们正是 \(\mathbb{Z}\) 在 \(\mathbb{Q}(\sqrt{D})\) 中的整闭包，如 15.3 节所述）。

在 \(D = -1\) 的特殊情形我们得到\textbf{高斯整数}环 \(\mathbb{Z}[i]\)，即 \(a, b\) 都是整数的复数 \(a+bi \in \mathbb{C}\). 这些数最初由 Gauss 约 1800 年引入，用以表述四次互反律，它讨论模素数的四次幂之间存在的优美关系。我们不久将看到这个环的代数结构在数论问题中的另一个有用应用。

定义\textbf{域范数} \(N : \mathbb{Q}(\sqrt{D}) \to \mathbb{Q}\) 为
\[
N(a + b\sqrt{D}) = (a+b\sqrt{D})(a-b\sqrt{D}) = a^2 - Db^2 \in \mathbb{Q},
\]
如前所述，若 \(a + b\sqrt{D} \neq 0\) 则它非零。这个范数给出域 \(\mathbb{Q}(\sqrt{D})\) 中“大小”的一种度量。例如当 \(D = -1\) 时 \(a+bi\) 的范数是 \(a^2+b^2\)，即这个复数看作复平面中向量时长度的平方。我们将在这个例子和后续例子中用域范数来建立环 \(\mathcal{O}\) 的许多性质。

容易验证 \(N\) 是乘性的，即对所有 \(\alpha, \beta \in \mathbb{Q}(\sqrt{D})\) 有 \(N(\alpha\beta) = N(\alpha)N(\beta)\). 在子环 \(\mathcal{O}\) 上也容易看出域范数由
\[
N(a + b\omega) = (a + b\omega)(a + b\bar\omega) = \begin{cases} a^2 - Db^2, & \text{若 } D \equiv 2, 3 \pmod 4, \\ a^2 + ab + \dfrac{1-D}{4}b^2, & \text{若 } D \equiv 1 \pmod 4 \end{cases}
\]
给出，其中
\[
\bar\omega = \begin{cases} \sqrt{D}, & \text{若 } D \equiv 2, 3 \pmod 4, \\ \dfrac{1-\sqrt{D}}{2}, & \text{若 } D \equiv 1 \pmod 4. \end{cases}
\]
由此推出对每个 \(\alpha \in \mathcal{O}\)，\(N(\alpha)\) 事实上是整数。

我们可以用这个范数刻画 \(\mathcal{O}\) 中的单位。若 \(\alpha \in \mathcal{O}\) 有域范数 \(N(\alpha) = \pm 1\)，上面的公式表明 \((a+b\omega)^{-1} = \pm(a + b\bar\omega)\)，它仍是 \(\mathcal{O}\) 的元素，故 \(\alpha\) 是 \(\mathcal{O}\) 中的单位。反之假设 \(\alpha\) 是 \(\mathcal{O}\) 中的单位，即对某个 \(\beta \in \mathcal{O}\) 有 \(\alpha\beta = 1\). 则域范数的乘性蕴含 \(N(\alpha)N(\beta) = N(\alpha\beta) = N(1) = 1\). 由于 \(N(\alpha)\) 和 \(N(\beta)\) 都是整数，两者都必为 \(\pm 1\). 因此
\begin{center}
元素 \(\alpha\) 是 \(\mathcal{O}\) 中的单位当且仅当 \(N(\alpha) = \pm 1\).
\end{center}
特别地，方程 \(x^2 - Dy^2 = \pm 1\) 的整数解（初等数论中称为 Pell 方程）的确定本质上等价于环 \(\mathcal{O}\) 中单位的确定。

当 \(D = -1\) 时，高斯整数 \(\mathbb{Z}[i]\) 中的单位是满足 \(a^2 + b^2 = \pm 1\)（\(a, b \in \mathbb{Z}\)）的元素 \(a+bi\)，故单位群由 \(\{\pm 1, \pm i\}\) 组成。当 \(D = -3\) 时，\(\mathbb{Z}[(1+\sqrt{-3})/2]\) 中的单位由满足 \(a^2 + ab + b^2 = \pm 1\) 的整数 \(a, b\) 确定，即 \((2a+b)^2 + 3b^2 = \pm 4\)，由此容易看出单位群是 6 阶群 \(\{\pm 1, \pm \rho, \pm \rho^2\}\)，其中 \(\rho = (-1+\sqrt{-3})/2\). 对任何其他 \(D < 0\)，类似地容易看出唯一的单位是 \(\{\pm 1\}\).

相反，当 \(D > 0\) 时可以证明单位群 \(\mathcal{O}^\times\) 总是无限的。例如，容易验证 \(1 + \sqrt{2}\) 是环 \(\mathcal{O} = \mathbb{Z}[\sqrt{2}]\) 中的单位（域范数为 \(-1\)），且 \(\{\pm(1+\sqrt{2})^n \mid n \in \mathbb{Z}\}\) 是无限的互不相同单位集合（事实上就是这个情形下的全部单位群，但这更难证明）。
\end{example}

\subsection*{习题}
设 \(R\) 是带单位元的环。

\begin{enumerate}
\item 证明 \((-1)^2 = 1\)（在 \(R\) 中）。

\item 证明：若 \(u\) 是 \(R\) 中的单位，则 \(-u\) 也是。

\item 设 \(R\) 是带单位元的环，\(S\) 是 \(R\) 的包含单位元的子环。证明：若 \(u\) 在 \(S\) 中是单位，则 \(u\) 在 \(R\) 中是单位。举例说明逆命题不成立。

\item 证明环的任意非空族子环的交也是子环。

\item 判断下列 (a)–(f) 中哪些是 \(\mathbb{Q}\) 的子环：

（a）所有分母为奇数（写成最简形式时）的有理数集合；

（b）所有分母为偶数（最简形式）的有理数集合；

（c）非负有理数集合；

（d）有理数的平方构成的集合；

（e）所有分子为奇数（最简形式）的有理数集合；

（f）所有分子为偶数（最简形式）的有理数集合。

\item 判断下列哪些是从闭区间 \([0,1]\) 到 \(\mathbb{R}\) 的所有函数环的子环：

（a）所有满足对所有 \(q \in \mathbb{Q} \cap [0,1]\) 有 \(f(q) = 0\) 的函数 \(f(x)\) 的集合；

（b）所有多项式函数的集合；

（c）所有只有有限多个零点的函数连同零函数的集合；

（d）所有有无穷多个零点的函数的集合；

（e）所有满足 \(\lim_{x \to 1^-} f(x) = 0\) 的函数 \(f\) 的集合；

（f）函数 \(\sin nx\) 和 \(\cos mx\)（\(m, n \in \{0, 1, 2, \ldots\}\)）的所有有理线性组合的集合。

\item 环 \(R\) 的\textbf{中心}是 \(\{z \in R \mid zr = rz \text{ 对所有 } r \in R\}\)（即所有与 \(R\) 的每个元素交换的元素构成的集合）。证明环的中心是包含单位元的子环。证明除环的中心是域。

\item 描述实 Hamilton 四元数的中心。证明 \(\{a + bi \mid a, b \in \mathbb{R}\}\) 是 \(\mathbb{H}\) 的子环，它是一个域但不含于 \(\mathbb{H}\) 的中心中。

\item 对固定的元素 \(a \in R\)，定义 \(C(a) = \{r \in R \mid ra = ar\}\). 证明 \(C(a)\) 是 \(R\) 的包含 \(a\) 的子环。证明 \(R\) 的中心是所有 \(a \in R\) 的子环 \(C(a)\) 的交。

\item 证明：若 \(D\) 是除环，则对所有 \(a \in D\)，\(C(a)\) 是除环（参见前一习题）。

\item 证明：若 \(R\) 是整环且对某个 \(x \in R\) 有 \(x^2 = 1\)，则 \(x = \pm 1\).

\item 证明任何域中包含单位元的子环都是整环。

\item \(R\) 中的元素 \(x\) 称为\textbf{幂零的}，若对某个 \(m \in \mathbb{Z}^+\) 有 \(x^m = 0\).

（a）证明：若对某些整数 \(a, b\) 有 \(n = ab\)，则 \(\overline{ab}\) 是 \(\mathbb{Z}/n\mathbb{Z}\) 中的幂零元素。

（b）若 \(a \in \mathbb{Z}\) 是整数，证明元素 \(\bar{a} \in \mathbb{Z}/n\mathbb{Z}\) 是幂零的当且仅当 \(n\) 的每个素因子也是 \(a\) 的因子。特别地，具体确定 \(\mathbb{Z}/72\mathbb{Z}\) 的幂零元素。

（c）设 \(R\) 是从非空集合 \(X\) 到域 \(F\) 的函数环。证明 \(R\) 不含非零幂零元素。

\item 设 \(x\) 是交换环 \(R\) 的幂零元素（参见前一习题）。

（a）证明 \(x\) 不是零就是零因子。

（b）证明对所有 \(r \in R\)，\(rx\) 是幂零的。

（c）证明 \(1 + x\) 是 \(R\) 中的单位。

（d）由此推出幂零元素与单位之和是单位。

\item 环 \(R\) 称为\textbf{布尔环}，若对所有 \(a \in R\) 有 \(a^2 = a\). 证明每个布尔环都交换。

\item 证明唯一的整环布尔环是 \(\mathbb{Z}/2\mathbb{Z}\).

\item 设 \(R\) 和 \(S\) 是环。证明直积 \(R \times S\) 在按分量的加法与乘法下是环。证明 \(R \times S\) 交换当且仅当 \(R\) 和 \(S\) 都交换。证明 \(R \times S\) 有单位元当且仅当 \(R\) 和 \(S\) 都有单位元。

\item 证明 \(\{(r, r) \mid r \in R\}\) 是 \(R \times R\) 的子环。

\item 设 \(I\) 是任意非空指标集，对每个 \(i \in I\)，\(R_i\) 是环。证明直积 \(\prod_{i \in I} R_i\) 在按分量的加法与乘法下是环。

\item 设 \(R\) 是整数序列 \((a_1, a_2, a_3, \ldots)\) 的集合，其中除有限多个 \(a_i\) 外都是 0（称为 \(\mathbb{Z}\) 的无限多个复制的直和）。证明 \(R\) 在按分量的加法与乘法下是无单位元的环。

\item 设 \(X\) 是任意非空集合，\(\mathcal{P}(X)\) 是 \(X\) 的所有子集的集合（\(X\) 的幂集）。在 \(\mathcal{P}(X)\) 上定义加法与乘法为
\[
A + B = (A - B) \cup (B - A), \qquad A \times B = A \cap B,
\]
即加法是对称差、乘法是交。

（a）证明 \(\mathcal{P}(X)\) 在这些运算下是环（\(\mathcal{P}(X)\) 及其子环常称为\textbf{集合环}）。

（b）证明这个环交换、有单位元且是布尔环。

\item 举一个无限布尔环的例子。

\item 设 \(D\) 是无平方因子的整数，\(\mathcal{O}\) 是二次域 \(\mathbb{Q}(\sqrt{D})\) 中的整数环。对任意正整数 \(f\)，证明集合 \(\mathcal{O}_f = \mathbb{Z}[f\omega] = \{a + bf\omega \mid a, b \in \mathbb{Z}\}\) 是 \(\mathcal{O}\) 的包含单位元的子环。证明 \([\mathcal{O} : \mathcal{O}_f] = f\)（作为加法阿贝尔群的指数）。反之证明：\(\mathcal{O}\) 的包含单位元且在 \(\mathcal{O}\) 中有有限指数 \(f\)（作为加法阿贝尔群）的子环等于 \(\mathcal{O}_f\).（环 \(\mathcal{O}_f\) 称为域 \(\mathbb{Q}(\sqrt{D})\) 中导子为 \(f\) 的序。整数环 \(\mathcal{O}\) 称为 \(\mathbb{Q}(\sqrt{D})\) 中的极大序。）

\item 对 \(D = 3, 5, 6, 7\)，通过在每个环中举出一个无穷（乘法）阶的单位，证明二次整数环 \(\mathcal{O}\) 的单位群 \(\mathcal{O}^\times\) 是无限的。

\item 设 \(I\) 是整 Hamilton 四元数环，定义
\[
N : I \to \mathbb{Z}, \quad N(a+bi+cj+dk) = a^2+b^2+c^2+d^2
\]
（映射 \(N\) 称为范数）。

（a）证明对所有 \(\alpha \in I\) 有 \(N(\alpha) = \alpha\bar\alpha\)，其中若 \(\alpha = a+bi+cj+dk\) 则 \(\bar\alpha = a-bi-cj-dk\).

（b）证明对所有 \(\alpha, \beta \in I\) 有 \(N(\alpha\beta) = N(\alpha)N(\beta)\).

（c）证明 \(I\) 的元素是单位当且仅当它的范数为 \(+1\). 证明 \(I^\times\) 同构于 8 阶四元数群。〔有理四元数环中非零元素 \(\alpha\) 的逆元是 \(\frac{\bar\alpha}{N(\alpha)}\).〕

\item 设 \(K\) 是域。\(K\) 上的\textbf{离散赋值}是满足下列条件的函数 \(v : K^\times \to \mathbb{Z}\)：

（i）\(v(ab) = v(a) + v(b)\)（即 \(v\) 是从 \(K\) 的非零元素的乘法群到 \(\mathbb{Z}\) 的同态）；

（ii）\(v\) 是满射；

（iii）对所有满足 \(x+y \neq 0\) 的 \(x, y \in K^\times\) 有 \(v(x+y) \geq \min\{v(x), v(y)\}\).

集合 \(R = \{x \in K^\times \mid v(x) \geq 0\} \cup \{0\}\) 称为 \(v\) 的\textbf{赋值环}。

（a）证明 \(R\) 是 \(K\) 的包含单位元的子环。（一般地，若存在某个域 \(K\) 和 \(K\) 上的某个离散赋值 \(v\) 使 \(R\) 是 \(v\) 的赋值环，则环 \(R\) 称为\textbf{离散赋值环}。）

（b）证明对每个非零元素 \(x \in K\)，\(x\) 或 \(x^{-1}\) 在 \(R\) 中。

（c）证明元素 \(x\) 是 \(R\) 的单位当且仅当 \(v(x) = 0\).

\item 离散赋值环的一个具体例子（参见前一习题）如下得到：当 \(p\) 是素数、\(K = \mathbb{Q}\) 且
\[
v\left(\frac{a}{b}\right) = c, \quad \text{其中 } \frac{a}{b} = p^c\,\frac{a'}{b'},\ p \nmid a',\ p \nmid b'
\]
时。证明相应的赋值环 \(R\) 是所有分母与 \(p\) 互素的有理数构成的环。描述这个赋值环的单位。

\item 设 \(R\) 是满足 \(1 \neq 0\) 的环。非零元素 \(a\) 称为 \(R\) 中的\textbf{左零因子}，若存在非零元素 \(x \in R\) 使 \(ax = 0\). 对称地，\(b \neq 0\) 称为\textbf{右零因子}，若存在非零 \(y \in R\) 使 \(yb = 0\)（故零因子是左零因子或右零因子的元素）。元素 \(u \in R\) 有\textbf{左逆}，若存在 \(s \in R\) 使 \(su = 1\). 对称地，\(v\) 有\textbf{右逆}，若对某个 \(t \in R\) 有 \(vt = 1\).

（a）证明 \(u\) 是单位当且仅当它既有右逆又有左逆（即 \(u\) 必须有双边逆）。

（b）证明：若 \(u\) 有右逆，则 \(u\) 不是右零因子。

（c）证明：若 \(u\) 有多于一个右逆，则 \(u\) 是左零因子。

（d）证明：若 \(R\) 是有限环，则每个有右逆的元素都是单位（即有双边逆）。

\item 设 \(A\) 是任意满足 \(1 \neq 0\) 的带单位元交换环。设 \(R\) 是 \(A\) 的加法群到自身的所有群同态的集合，加法由函数的逐点加法定义、乘法由函数复合定义。证明这些运算使 \(R\) 成为带单位元的环。证明 \(R\) 的单位是 \(A\) 的群自同构（参见第 1.6 节习题 20）。

\item 设 \(A = \mathbb{Z} \times \mathbb{Z} \times \mathbb{Z} \times \cdots\) 是由正整数指标标记的 \(\mathbb{Z}\) 的复制的直积（故 \(A\) 在按分量的加法与乘法下是环），设 \(R\) 是前一习题中描述的 \(A\) 到自身的所有群同态构成的环。设 \(\varphi\) 是 \(R\) 中由 \(\varphi(a_1, a_2, a_3, \ldots) = (a_2, a_3, \ldots)\) 定义的元素。设 \(\psi\) 是 \(R\) 中由 \(\psi(a_1, a_2, a_3, \ldots) = (0, a_1, a_2, a_3, \ldots)\) 定义的元素。

（a）证明 \(\varphi\psi\) 是 \(R\) 的单位元但 \(\psi\varphi\) 不是 \(R\) 的单位元（即 \(\psi\) 是 \(\varphi\) 的右逆但不是左逆）。

（b）举出 \(\varphi\) 的无穷多个右逆。

（c）求 \(R\) 中一个非零元素 \(\pi\) 使 \(\varphi\pi = 0\) 但 \(\pi\varphi \neq 0\).

（d）证明不存在非零元素 \(\lambda \in R\) 使 \(\lambda\varphi = 0\)（即 \(\varphi\) 是左零因子但不是右零因子）。
\end{enumerate}

\section{例子：多项式环、矩阵环与群环}

本节我们引入三类重要的环：多项式环、矩阵环和群环。我们将在全书过程中看到这三类环常常相互联系。例如，我们将在第 VI 部分看到群 \(G\) 在复数上的群环是 \(\mathbb{C}\) 上矩阵环的直积。

这些环除了本身有趣外还有许多重要应用。在第 III 部分我们将用多项式环证明矩阵的一些分类定理，它们特别地确定一个矩阵何时相似于对角矩阵。在第 VI 部分我们将用群环研究群作用并证明一些额外的分类定理。

\subsection*{多项式环}

固定带单位元的交换环 \(R\). 设 \(x\) 是未定元。形式
\[
\sum_{i=0}^{n} a_i x^i = a_nx^n + a_{n-1}x^{n-1} + \cdots + a_1x + a_0
\]
（其中 \(n \geq 0\) 且每个 \(a_i \in R\)）称为 \(x\) 的、系数 \(a_i\) 在 \(R\) 中的\textbf{多项式}。若 \(a_n \neq 0\)，则称这个多项式的\textbf{次数}为 \(n\)，\(a_nx^n\) 称为\textbf{首项}，\(a_n\) 称为\textbf{首项系数}（零多项式的首项系数取为 0）。若 \(a_n = 1\)，则称多项式是\textbf{首一的}。所有这样的多项式的集合称为系数在 \(R\) 中的、变量 \(x\) 的\textbf{多项式环}，记作 \(R[x]\).

使 \(R[x]\) 成为环的加法与乘法运算与初等代数中熟悉的运算相同：加法是“逐分量”的
\[
\begin{aligned}
&(a_nx^n + a_{n-1}x^{n-1} + \cdots + a_1x + a_0) + (b_nx^n + b_{n-1}x^{n-1} + \cdots + b_1x + b_0) \\
&= (a_n+b_n)x^n + (a_{n-1}+b_{n-1})x^{n-1} + \cdots + (a_1+b_1)x + (a_0+b_0)
\end{aligned}
\]
（这里为了定义不同次数的多项式的加法，\(a_n\) 或 \(b_n\) 可以为零）。乘法首先对只有一个非零项的多项式定义 \((ax^i)(bx^j) = abx^{i+j}\)，然后通过分配律推广到所有多项式（通常称为“展开并合并同类项”）：
\[
(a_0 + a_1x + a_2x^2 + \cdots) \times (b_0 + b_1x + b_2x^2 + \cdots)
\]
\[
= a_0b_0 + (a_0b_1 + a_1b_0)x + (a_0b_2 + a_1b_1 + a_2b_0)x^2 + \cdots
\]
（一般地，乘积中 \(x^k\) 的系数是 \(\sum_{i=0}^{k} a_ib_{k-i}\)）。这些运算有意义，因为 \(R\) 是环，故系数的和与积有定义。容易验证 \(R[x]\) 在这些加法与乘法定义下确实是环。

环 \(R\) 以常多项式出现在 \(R[x]\) 中。注意由乘法的定义，\(R[x]\) 是带单位元（\(R\) 的单位元 1）的交换环。

上面假设系数环 \(R\) 是交换环，因为这是我们主要关心的情形；但注意即使 \(R\) 不交换或没有单位元，上面 \(R[x]\) 中加法与乘法的定义仍然有效。若系数环 \(R\) 是整数 \(\mathbb{Z}\)（分别为有理数 \(\mathbb{Q}\)），则多项式环 \(\mathbb{Z}[x]\)（分别为 \(\mathbb{Q}[x]\)）是初等代数中熟悉的整（有理）系数多项式环。

另一个例子是系数在 \(\mathbb{Z}/3\mathbb{Z}\) 中的 \(x\) 的多项式环 \(\mathbb{Z}/3\mathbb{Z}[x]\). 这个环由 \(x\) 的非负幂组成，系数为 0, 1, 2，对系数的计算按模 3 进行。例如，若
\[
p(x) = x^2 + 2x + 1, \qquad q(x) = x^3 + x + 2,
\]
则
\[
p(x) + q(x) = x^3 + x^2
\]
且
\[
p(x)q(x) = x^5 + 2x^4 + 2x^3 + x^2 + 2x + 2.
\]

取系数的环对多项式的行为有很大影响。例如，多项式 \(x^2 + 1\) 在多项式环 \(\mathbb{Z}[x]\) 中不是完全平方，但在多项式环 \(\mathbb{Z}/2\mathbb{Z}[x]\) 中是，因为在这个环中 \((x+1)^2 = x^2 + 2x + 1 = x^2 + 1\).

\begin{proposition}\label{pro:7.4}
设 \(R\) 是整环，\(p(x), q(x)\) 是 \(R[x]\) 的非零元素。则

\begin{enumerate}[label=(\arabic*)]
\item \(\deg p(x)q(x) = \deg p(x) + \deg q(x)\)；
\item \(R[x]\) 的单位就是 \(R\) 的单位；
\item \(R[x]\) 是整环。
\end{enumerate}
\end{proposition}

\begin{proof}
若 \(R\) 没有零因子，则 \(R[x]\) 也没有；若 \(p(x)\) 和 \(q(x)\) 的首项分别为 \(a_nx^n\) 和 \(b_mx^m\)，则 \(p(x)q(x)\) 的首项是 \(a_nb_mx^{n+m}\)，且 \(a_nb_m \neq 0\). 这证明了（3）也验证了（1）。若 \(p(x)\) 是单位，即 \(R[x]\) 中有 \(p(x)q(x) = 1\)，则 \(\deg p(x) + \deg q(x) = 0\)，故 \(p(x)\) 和 \(q(x)\) 都是 \(R\) 的元素，从而由于它们的乘积为 1 而是 \(R\) 中的单位。这证明了（2）。
\end{proof}

若环 \(R\) 有零因子，则 \(R[x]\) 也有，因为 \(R \subseteq R[x]\). 此外，若 \(f(x)\) 是 \(R[x]\) 中的零因子（即对某个非零 \(g(x) \in R[x]\) 有 \(f(x)g(x) = 0\)），则事实上对某个非零 \(c \in R\) 有 \(cf(x) = 0\)（参见习题 2）。

若 \(S\) 是 \(R\) 的子环，则 \(S[x]\) 是 \(R[x]\) 的子环。例如 \(\mathbb{Z}[x]\) 是 \(\mathbb{Q}[x]\) 的子环。\(R[x]\) 的子环的其他例子是所有 \(x^2\) 的多项式（即只出现 \(x\) 的偶次幂）的集合，以及所有常数项为零的多项式的集合（后一个子环没有单位元）。

多项式环，特别是域上的多项式环，将在第 9 章广泛研究。

\subsection*{矩阵环}

固定任意环 \(R\)，设 \(n\) 是正整数。设 \(M_n(R)\) 是所有元素取自 \(R\) 的 \(n \times n\) 矩阵的集合。\(M_n(R)\) 的元素 \((a_{ij})\) 是 \(R\) 的元素的 \(n \times n\) 方阵，其第 \(i\) 行第 \(j\) 列的元素是 \(a_{ij} \in R\). 矩阵集合在实数矩阵通常的加法与乘法规则下成为环。加法按分量：矩阵 \((a_{ij}) + (b_{ij})\) 的 \(i, j\) 元素是 \(a_{ij} + b_{ij}\). 矩阵乘积 \((a_{ij}) \times (b_{ij})\) 的 \(i, j\) 元素是 \(\sum_{k=1}^{n} a_{ik}b_{kj}\)（注意这些矩阵必须是方阵，才能定义任意两个元素的乘法）。容易直接验证这些运算使 \(M_n(R)\) 成为环。当 \(R\) 是域时，我们将在第 III 部分用较少计算的方法证明 \(M_n(R)\) 是环。

注意若 \(R\) 是任意非平凡环（即使是交换环）且 \(n \geq 2\)，则 \(M_n(R)\) 不交换：若 \(ab \neq 0\)（在 \(R\) 中），设 \(A\) 是 (1,1) 位置为 \(a\)、其余为零的矩阵，\(B\) 是 (1,2) 位置为 \(b\)、其余为零的矩阵；则 \(AB\) 的 (1,2) 位置是（非零的）\(ab\)，而 \(BA\) 是零矩阵。这两个矩阵也表明当 \(n \geq 2\) 时对任意非零环 \(R\)，\(M_n(R)\) 有零因子。

\(M_n(R)\) 的元素 \((a_{ij})\) 称为\textbf{纯量矩阵}，若对某个 \(a \in R\) 有 \(a_{ii} = a\)（对所有 \(i \in \{1, \ldots, n\}\)）且所有 \(i \neq j\) 有 \(a_{ij} = 0\)（即所有对角元都等于 \(a\)、所有非对角元为零）。纯量矩阵的集合是 \(M_n(R)\) 的子环。这个子环是 \(R\) 的一个复制（即“同构”于 \(R\)）：若矩阵 \(A\) 的主对角元为 \(a\)、矩阵 \(B\) 的主对角元为 \(b\)，则矩阵 \(A+B\) 的对角元为 \(a+b\)、\(AB\) 的对角元为 \(ab\)（其余元素为零）。若 \(R\) 交换，则纯量矩阵与 \(M_n(R)\) 的所有元素交换。若 \(R\) 有 1，则主对角线上全为 1 的纯量矩阵（即 \(n \times n\) 单位矩阵）是 \(M_n(R)\) 的 1. 此时 \(M_n(R)\) 中的单位是可逆 \(n \times n\) 矩阵，单位群记作 \(GL_n(R)\)，即 \(R\) 上 \(n\) 次\textbf{一般线性群}。

若 \(S\) 是 \(R\) 的子环，则 \(M_n(S)\) 是 \(M_n(R)\) 的子环。例如 \(M_n(\mathbb{Z})\) 是 \(M_n(\mathbb{Q})\) 的子环，而 \(M_n(2\mathbb{Z})\) 是这两者的子环。\(M_n(R)\) 的子环的另一个例子是上三角矩阵的集合：\(\{(a_{ij}) \mid \text{只要 } p > q \text{ 就有 } a_{pq} = 0\}\)（即所有主对角线以下元素为零的矩阵的集合）——容易验证上三角矩阵的和与积仍是上三角的。

\subsection*{群环}

固定带单位元 \(1 \neq 0\) 的交换环 \(R\)，设 \(G = \{g_1, g_2, \ldots, g_n\}\) 是任意有限群，群运算写成乘法。定义 \(G\) 的、系数在 \(R\) 中的\textbf{群环} \(RG\) 为所有形式和
\[
a_1g_1 + a_2g_2 + \cdots + a_ng_n, \quad a_i \in R,\ 1 \leq i \leq n
\]
的集合。若 \(g_1\) 是 \(G\) 的单位元，我们把 \(a_1g_1\) 简写为 \(a_1\). 类似地，对 \(g \in G\) 把元素 \(1g\) 简写为 \(g\).

加法“按分量”定义
\[
(a_1g_1 + a_2g_2 + \cdots + a_ng_n) + (b_1g_1 + b_2g_2 + \cdots + b_ng_n) = (a_1+b_1)g_1 + (a_2+b_2)g_2 + \cdots + (a_n+b_n)g_n.
\]
乘法首先定义 \((ag_i)(bg_j) = (ab)g_k\)，其中乘积 \(ab\) 在 \(R\) 中取，\(g_ig_j = g_k\) 是群 \(G\) 中的乘积。然后通过分配律把这个乘积推广到所有形式和，于是乘积 \((a_1g_1 + \cdots + a_ng_n) \times (b_1g_1 + \cdots + b_ng_n)\) 中 \(g_k\) 的系数是 \(\sum_{g_ig_j = g_k} a_ib_j\). 容易验证这些运算使 \(RG\) 成为环（同样不需要 \(R\) 交换）。乘法的结合性由群 \(G\) 中群运算的结合性推出。环 \(RG\) 交换当且仅当 \(G\) 是交换群。

\begin{example}
设 \(G = D_8\) 是 8 阶二面体群，通常生成元为 \(r, s\)（\(r^4 = s^2 = 1\) 且 \(rs = sr^{-1}\)），\(R = \mathbb{Z}\). 元素 \(\alpha = r + r^2 - 2s\) 和 \(\beta = -3r^2 + rs\) 是 \(\mathbb{Z}D_8\) 的典型成员。它们的和与积为
\[
\alpha + \beta = r - 2r^2 - 2s + rs,
\]
\[
\begin{aligned}
\alpha\beta &= (r + r^2 - 2s)(-3r^2 + rs) \\
&= r(-3r^2 + rs) + r^2(-3r^2 + rs) - 2s(-3r^2 + rs) \\
&= -3r^3 + r^2s - 3 + r^3s + 6r^2s - 2r^3 = -3 - 5r^3 + 7r^2s + r^3s.
\end{aligned}
\]
环 \(R\) 以“常”形式和出现在 \(RG\) 中，即 \(G\) 的单位元的 \(R\)-倍数（注意限制到这些元素的 \(RG\) 中加法与乘法的定义就是 \(R\) 中的加法与乘法）。这些 \(R\) 的元素与 \(RG\) 的所有元素交换。\(R\) 的单位元是 \(RG\) 的单位元。

群 \(G\) 也出现在 \(RG\) 中（例如元素 \(g_i\) 表现为 \(1g_i\)——上面的 \(r, s \in D_8\) 也是群环 \(\mathbb{Z}D_8\) 的元素）——限制到 \(G\) 的 \(RG\) 中乘法就是群运算。特别地，\(G\) 的每个元素在群环 \(RG\) 中都有乘法逆元（即它在 \(G\) 中的逆元）。这表明 \(G\) 是 \(RG\) 的单位群的子群。

若 \(|G| > 1\)，则 \(RG\) 总有零因子。例如设 \(g\) 是 \(G\) 的任意 \(m > 1\) 阶元素。则
\[
(1-g)(1 + g + \cdots + g^{m-1}) = 1 - g^m = 1 - 1 = 0,
\]
故 \(1-g\) 是零因子（注意按 \(RG\) 的定义，上面乘积中的两个形式和都不为零）。

若 \(S\) 是 \(R\) 的子环，则 \(SG\) 是 \(RG\) 的子环。例如 \(\mathbb{Z}G\)（称为 \(G\) 的\textbf{整群环}）是 \(\mathbb{Q}G\)（\(G\) 的\textbf{有理群环}）的子环。此外，若 \(H\) 是 \(G\) 的子群，则 \(RH\) 是 \(RG\) 的子环。\(RG\) 中系数之和为零的所有元素构成的集合是子环（无单位元）。若 \(|G| > 1\)，则“常数项”（即 \(G\) 的单位元的系数）为零的元素集合不是子环（它不对乘法封闭）。

注意群环 \(\mathbb{R}Q_8\) 与 Hamilton 四元数 \(\mathbb{H}\) 不是同一个环，尽管后者（在乘法下）含有四元数群 \(Q_8\) 的一个复制。一个区别是 \(Q_8\) 中唯一的 2 阶元素（通常记作 \(-1\)）不是 \(\mathbb{R}Q_8\) 中 1 的加法逆元。换句话说，若我们暂时把群 \(Q_8\) 的单位元记作 \(g_1\)、唯一的 2 阶元素记作 \(g_2\)，则 \(g_1 + g_2\) 在 \(\mathbb{R}Q_8\) 中不为零，而 \(1 + (-1)\) 在 \(\mathbb{H}\) 中为零。此外，如上所述，群环 \(\mathbb{R}Q_8\) 含零因子，故不是除环。
\end{example}

域上的群环将在第 18 章广泛研究。

\subsection*{习题}
设 \(R\) 是带 1 的交换环。

\begin{enumerate}
\item 设 \(p(x) = 2x^3 - 3x^2 + 4x - 5\)，\(q(x) = 7x^3 + 33x - 4\). 在 (a)、(b)、(c) 中各计算 \(p(x)+q(x)\) 和 \(p(x)q(x)\)，假设给定多项式的系数取自指定的环（其中整数系数在 (b) 和 (c) 中取模）：

（a）\(R = \mathbb{Z}\)；（b）\(R = \mathbb{Z}/2\mathbb{Z}\)；（c）\(R = \mathbb{Z}/3\mathbb{Z}\).

\item 设 \(p(x) = a_nx^n + a_{n-1}x^{n-1} + \cdots + a_1x + a_0\) 是多项式环 \(R[x]\) 的元素。证明 \(p(x)\) 是 \(R[x]\) 中的零因子当且仅当存在非零 \(b \in R\) 使 \(bp(x) = 0\).〔设 \(g(x) = b_mx^m + b_{m-1}x^{m-1} + \cdots + b_0\) 是使 \(g(x)p(x) = 0\) 的非零最小次数多项式。证明 \(b_ma_n = 0\)，故 \(a_ng(x)\) 是次数小于 \(m\) 且乘以 \(p(x)\) 也得 0 的多项式。推出 \(a_ng(x) = 0\). 对 \(i\) 作归纳用类似的论证证明 \(a_{n-i}g(x) = 0\)（\(i = 0, 1, \ldots, n\)），并证明这意味着 \(b_mp(x) = 0\).〕

\item 定义 \(R\) 上未定元 \(x\) 的\textbf{形式幂级数}集合 \(R[[x]]\) 为所有形式和
\[
\sum_{n=0}^{\infty} a_nx^n = a_0 + a_1x + a_2x^2 + a_3x^3 + \cdots.
\]
按与实或复系数幂级数相同的方式定义幂级数的加法与乘法，即把多项式加法与乘法推广到幂级数，好像它们是“无限次多项式”：
\[
\sum_{n=0}^{\infty} a_nx^n + \sum_{n=0}^{\infty} b_nx^n = \sum_{n=0}^{\infty}(a_n+b_n)x^n,
\]
\[
\sum_{n=0}^{\infty} a_nx^n \times \sum_{n=0}^{\infty} b_nx^n = \sum_{n=0}^{\infty}\left(\sum_{k=0}^{n} a_kb_{n-k}\right)x^n.
\]
（这里用“形式”一词表明不考虑收敛性，故形式幂级数不必表示 \(R\) 上的函数。）

（a）证明 \(R[[x]]\) 是带 1 的交换环。

（b）证明 \(1-x\) 是 \(R[[x]]\) 中的单位，逆元为 \(1+x+x^2+\cdots\).

（c）证明 \(\sum_{n} a_nx^n\) 是 \(R[[x]]\) 中的单位当且仅当 \(a_0\) 是 \(R\) 中的单位。

\item 证明若 \(R\) 是整环，则形式幂级数环 \(R[[x]]\) 也是整环。

\item 设 \(F\) 是域，定义系数取自 \(F\) 的\textbf{形式洛朗级数}环 \(F((x))\) 为
\[
F((x)) = \left\{\sum_{n \geq N} a_nx^n \ \middle|\ a_n \in F,\ N \in \mathbb{Z}\right\}.
\]
（\(F((x))\) 的每个元素都是 \(x\) 的幂级数加上 \(1/x\) 的多项式，即 \(F((x))\) 的每个元素只有有限多个负幂项。）

（a）证明 \(F((x))\) 是域。

（b）定义映射
\[
v : F((x))^\times \to \mathbb{Z}, \quad v\left(\sum_{n \geq N} a_nx^n\right) = N,
\]
其中 \(a_N\) 是级数的第一个非零系数（即 \(N\) 是“级数在 0 处的零点或极点的阶”）。证明 \(v\) 是 \(F((x))\) 上的离散赋值，其离散赋值环是形式幂级数环 \(F[[x]]\)（参见第 1 节习题 26）。

\item 设 \(S\) 是带单位元 \(1 \neq 0\) 的环。设 \(E_{ij}\) 表示 \(M_n(S)\) 中 \((i, j)\) 位置为 1、其余元素全为 0 的 \(n \times n\) 矩阵（\(1 \leq i, j \leq n\)），设 \(A\) 是 \(M_n(S)\) 中元素为 \(a_{ij}\) 的任意矩阵。

（a）证明 \(E_{ij}A\) 是第 \(i\) 行等于 \(A\) 的第 \(j\) 行、其余行全为零的矩阵。

（b）证明 \(AE_{ij}\) 是第 \(j\) 列等于 \(A\) 的第 \(i\) 列、其余列全为零的矩阵。

（c）由此推出 \(E_{pq}AE_{rs}\) 是 \((p, s)\) 位置为 \(a_{qr}\)、其余元素全为零的矩阵。

\item 证明环 \(M_n(R)\) 的中心是纯量矩阵的集合（参见第 1 节习题 7）。〔用前一习题。〕

\item 设 \(S\) 是任意环，\(n \geq 2\) 是整数。证明：若 \(A\) 是 \(M_n(S)\) 中的任意严格上三角矩阵，则 \(A^n = 0\)（严格上三角矩阵是主对角线上及其以下元素全为零的矩阵）。

\item 设 \(\alpha = r + r^2 - 2s\)、\(\beta = -3r^2 + rs\) 是本节描述的整群环 \(\mathbb{Z}D_8\) 的两个元素。计算 \(\mathbb{Z}D_8\) 的下列元素：

（a）\(\beta\alpha\)；（b）\(\alpha^2\)；（c）\(\alpha\beta - \beta\alpha\)；（d）\(\beta\alpha\beta\).

\item 考虑整群环 \(\mathbb{Z}S_3\) 的下列元素：
\[
\alpha = 3(1\,2) - 5(2\,3) + 14(1\,2\,3), \qquad \beta = 6(1) + 2(2\,3) - 7(1\,3\,2)
\]
（其中 \((1)\) 是 \(S_3\) 的单位元）。计算下列元素：

（a）\(\alpha+\beta\)；（b）\(2\alpha - 3\beta\)；（c）\(\alpha\beta\)；（d）\(\beta\alpha\)；（e）\(\alpha^2\).

\item 假设 \(\alpha\) 和 \(\beta\) 的系数在 \(\mathbb{Z}/3\mathbb{Z}\) 中（即 \(\alpha, \beta \in \mathbb{Z}/3\mathbb{Z}S_3\)），重复前一习题。

\item 设 \(G = \{g_1, \ldots, g_n\}\) 是有限群。证明群环 \(RG\) 中的元素 \(N = g_1 + g_2 + \cdots + g_n\) 在 \(RG\) 的中心中（参见第 1 节习题 7）。

\item 设 \(K = \{k_1, \ldots, k_m\}\) 是有限群 \(G\) 中的一个共轭类。

（a）证明元素 \(\hat{K} = k_1 + \cdots + k_m\) 在群环 \(RG\) 的中心中（参见第 1 节习题 7）。〔验证对所有 \(g \in G\) 有 \(g^{-1}Kg = K\).〕

（b）设 \(K_1, \ldots, K_r\) 是 \(G\) 的共轭类，对每个 \(K_i\) 设 \(\hat{K}_i\) 是 \(RG\) 中 \(K_i\) 的成员之和。证明元素 \(\alpha \in RG\) 在 \(RG\) 的中心中当且仅当对某些 \(a_1, a_2, \ldots, a_r \in R\) 有 \(\alpha = a_1\hat{K}_1 + a_2\hat{K}_2 + \cdots + a_r\hat{K}_r\).
\end{enumerate}

\section{环同态与商环}

环同态是从一个环到另一个环的、保持加法与乘法结构的映射：

\begin{definition}
设 \(R\) 和 \(S\) 是环。

（1）\textbf{环同态}是满足下列条件的映射 \(\varphi : R \to S\)：

（i）对所有 \(a, b \in R\) 有 \(\varphi(a+b) = \varphi(a)+\varphi(b)\)（故 \(\varphi\) 是加法群上的群同态），且

（ii）对所有 \(a, b \in R\) 有 \(\varphi(ab) = \varphi(a)\varphi(b)\).

（2）环同态 \(\varphi\) 的\textbf{核}，记作 \(\ker \varphi\)，是 \(R\) 中映到 \(S\) 中 0 的元素集合（即把 \(\varphi\) 看作加法群同态时的核）。

（3）双射的环同态称为\textbf{同构}。
\end{definition}

若上下文清楚，我们简单地用“同态”而不是“环同态”。类似地，若 \(A\) 和 \(B\) 是环，除非另有说明，\(A \cong B\) 总是指环的同构。

\begin{example}
（1）由把偶数映到 0、奇数映到 1 定义的映射 \(\varphi : \mathbb{Z} \to \mathbb{Z}/2\mathbb{Z}\) 是环同态。这个映射是加性的，因为两个偶数或两个奇数之和是偶数，一个偶数与一个奇数之和是奇数。这个映射是乘性的，因为两个奇数之积是奇数，一个偶数与任何整数之积是偶数。\(\varphi\) 的核（\(\varphi\) 在 \(0 \in \mathbb{Z}/2\mathbb{Z}\) 上方的纤维）是偶数集合。\(\varphi\) 在 \(1 \in \mathbb{Z}/2\mathbb{Z}\) 上方的纤维是奇数集合。

（2）对 \(n \in \mathbb{Z}\)，由 \(\varphi_n(x) = nx\) 定义的映射 \(\varphi_n : \mathbb{Z} \to \mathbb{Z}\) 一般不都是环同态，因为 \(\varphi_n(xy) = nxy\)，而 \(\varphi_n(x)\varphi_n(y) = nxny = n^2xy\). 故 \(\varphi_n\) 是环同态只有当 \(n^2 = n\)，即 \(n = 0, 1\) 时。然而注意 \(\varphi_n\) 总是加法群上的群同态。因此处理环时应小心，要确保两个环运算都被保持。注意 \(\varphi_0\) 是零同态，\(\varphi_1\) 是恒等同态。

（3）设 \(\varphi : \mathbb{Q}[x] \to \mathbb{Q}\) 是由 \(\varphi(p(x)) = p(0)\) 定义的从有理系数 \(x\) 的多项式环到有理数的映射（即把多项式映到它的常数项）。则 \(\varphi\) 是环同态，因为两个多项式之和的常数项是它们常数项之和，两个多项式之积的常数项是它们常数项之积。\(a \in \mathbb{Q}\) 上方的纤维由以 \(a\) 为常数项的多项式集合组成。特别地，\(\varphi\) 的核由常数项为 0 的多项式组成。
\end{example}

\begin{proposition}\label{pro:7.5}
设 \(R\) 和 \(S\) 是环，\(\varphi : R \to S\) 是同态。

\begin{enumerate}[label=(\arabic*)]
\item \(\varphi\) 的像是 \(S\) 的子环。
\item \(\varphi\) 的核是 \(R\) 的子环。此外，若 \(a \in \ker \varphi\)，则对所有 \(r \in R\) 有 \(ra, ar \in \ker \varphi\)，即 \(\ker \varphi\) 对乘以 \(R\) 中元素封闭。
\end{enumerate}
\end{proposition}

\begin{proof}
（1）若 \(s_1, s_2 \in \operatorname{im}\varphi\)，则对某些 \(r_1, r_2 \in R\) 有 \(s_1 = \varphi(r_1)\)、\(s_2 = \varphi(r_2)\). 则 \(\varphi(r_1 - r_2) = s_1 - s_2\) 且 \(\varphi(r_1r_2) = s_1s_2\). 这表明 \(s_1 - s_2, s_1s_2 \in \operatorname{im}\varphi\)，故 \(\varphi\) 的像对减法与乘法封闭，从而是 \(S\) 的子环。

（2）若 \(\alpha, \beta \in \ker\varphi\)，则 \(\varphi(\alpha) = \varphi(\beta) = 0\). 于是 \(\varphi(\alpha - \beta) = 0\) 且 \(\varphi(\alpha\beta) = 0\)，故 \(\ker\varphi\) 对减法与乘法封闭，是 \(R\) 的子环。类似地，对任意 \(r \in R\) 有 \(\varphi(r\alpha) = \varphi(r)\varphi(\alpha) = \varphi(r)0 = 0\)，且 \(\varphi(\alpha r) = \varphi(\alpha)\varphi(r) = 0\varphi(r) = 0\)，故 \(r\alpha, \alpha r \in \ker\varphi\).
\end{proof}

在群同态 \(\varphi\) 的情形我们看到同态的纤维具有自然同构于 \(\varphi\) 的像的群结构，这引出了对正规子群取商群的概念。对环的同态有类似的结果。

设 \(\varphi : R \to S\) 是核为 \(I\) 的环同态。由于 \(R\) 和 \(S\) 特别是加法阿贝尔群，\(\varphi\) 特别是阿贝尔群的同态，\(\varphi\) 的纤维是核 \(I\) 的加法陪集 \(r + I\)（更确切地说，若 \(r\) 是 \(R\) 中映到 \(a \in S\) 的任意元素，\(\varphi(r) = a\)，则 \(\varphi\) 在 \(a\) 上方的纤维是核 \(I\) 的陪集 \(r+I\)）。这些纤维具有自然同构于 \(\varphi\) 的像的环结构：若 \(X\) 是 \(a \in S\) 上方的纤维、\(Y\) 是 \(b \in S\) 上方的纤维，则 \(X+Y\) 是 \(a+b\) 上方的纤维，\(XY\) 是 \(ab\) 上方的纤维。用核 \(I\) 的陪集表示，这个加法与乘法是
\[
(r+I) + (s+I) = (r+s)+I, \tag{7.1}
\]
\[
(r+I) \times (s+I) = (rs)+I. \tag{7.2}
\]
与群的情形一样，验证这些运算在核 \(I\) 的陪集集合上定义环结构最终依赖于 \(S\) 的相应环性质。这个陪集环称为 \(R\) 对 \(I = \ker\varphi\) 的\textbf{商环}，记作 \(R/I\). 注意环 \(R/I\) 的加法结构就是加法阿贝尔群 \(R\) 对（必然正规的）子群 \(I\) 的加法商群。当 \(I\) 是某个同态 \(\varphi\) 的核时，这个加法阿贝尔商群还具有由 (7.2) 定义的乘法结构，使 \(R/I\) 成为环。

与群的情形一样，我们也可以考虑能否用（1）和（2）在 \(R\) 的任意子群 \(I\) 的陪集集合上定义环结构。注意由于 \(R\) 是阿贝尔加法群，子群 \(I\) 必然正规，故 \(I\) 的陪集的商 \(R/I\) 自动是加法阿贝尔群。问题在于这个商群是否还具有由 \(R\) 中乘法通过（2）诱导的乘法结构。一般答案是否定的（正如试图对群的任意子群取商的答案也是否定的），这引出了 \(R\) 中\textbf{理想}的概念（环中正规子群的类比）。然后我们将看到 \(R\) 的理想恰好是 \(R\) 的环同态的核（环中“正规子群是群同态的核”这一定理的类比）。

设 \(I\) 是加法群 \(R\) 的任意子群。我们考虑何时（2）中陪集的乘法良定义并使加法阿贝尔群 \(R/I\) 成为环。（2）中乘法良定义这一陈述是说乘法与所选取的特定代表元 \(r\) 和 \(s\) 无关，即若改用代表元 \(r+\alpha\) 和 \(s+\beta\)（\(\alpha, \beta \in I\)），右端得到同一个陪集。换句话说，我们必须有
\[
(r+\alpha)(s+\beta) + I = rs + I
\]
对所有 \(r, s \in R\) 和所有 \(\alpha, \beta \in I\) 成立。

取 \(r = s = 0\)，我们看到 \(I\) 必须对乘法封闭，即 \(I\) 必须是 \(R\) 的子环。

其次，取 \(s = 0\) 并让 \(r\) 任意，我们看到必须对每个 \(r \in R\) 和每个 \(\beta \in I\) 有 \(r\beta \in I\)，即 \(I\) 必须对左乘 \(R\) 中元素封闭。取 \(r = 0\) 并让 \(s\) 任意，我们类似地看到 \(I\) 必须对右乘 \(R\) 中元素封闭。

反之，若 \(I\) 对左乘和右乘 \(R\) 中元素都封闭，则上述关系对所有 \(\alpha, \beta \in I\) 成立。因此这是（2）中乘法良定义的充要条件。

最后，若由（2）定义的陪集乘法良定义，则这个乘法使加法商群 \(R/I\) 成为环。商中的每条环公理由 \(R\) 中相应公理直接推出。例如，分配律之一验证如下：
\[
\begin{aligned}
(r+I)\left[(s+I)+(t+I)\right] &= (r+I)\left[(s+t)+I\right] = r(s+t)+I = (rs+rt)+I \\
&= (rs+I) + (rt+I) = [(r+I)(s+I)] + [(r+I)(t+I)].
\end{aligned}
\]
这表明环 \(R\) 对子群 \(I\) 的商 \(R/I\) 有自然的环结构，当且仅当 \(I\) 也对左乘和右乘 \(R\) 中元素封闭（特别地 \(I\) 必须是 \(R\) 的子环，因为它对乘法封闭）。如上所述，这样的子环 \(I\) 称为 \(R\) 的\textbf{理想}：

\begin{definition}
设 \(R\) 是环，\(I\) 是 \(R\) 的子集，\(r \in R\).

（1）\(rI = \{ra \mid a \in I\}\)，\(Ir = \{ar \mid a \in I\}\).

（2）\(R\) 的子集 \(I\) 是 \(R\) 的\textbf{左理想}，若

（i）\(I\) 是 \(R\) 的子环，且

（ii）\(I\) 对左乘 \(R\) 中元素封闭，即对所有 \(r \in R\) 有 \(rI \subseteq I\).

类似地，\(I\) 是\textbf{右理想}，若（i）成立且把（ii）换成

（ii）\('\) \(I\) 对右乘 \(R\) 中元素封闭，即对所有 \(r \in R\) 有 \(Ir \subseteq I\).

（3）既是左理想又是右理想的子集 \(I\) 称为 \(R\) 的\textbf{理想}（或为强调起见称为\textbf{双边理想}）。
\end{definition}

对交换环，左理想、右理想和双边理想的概念重合。我们强调，要证明环 \(R\) 的子集 \(I\) 是理想，必须证明 \(I\) 非空、对减法封闭、且对乘以 \(R\) 的所有元素（而不只是 \(I\) 的元素）封闭。若 \(R\) 有 1，则 \((-1)a = -a\)，故此时 \(I\) 非空、对加法封闭且对乘以 \(R\) 的所有元素封闭就足以保证它是理想。

还要注意命题 \ref{pro:7.5} 的最后一部分证明了任何环同态的核都是理想。

我们把前面的讨论总结为下面的命题。

\begin{proposition}\label{pro:7.6}
设 \(R\) 是环，\(I\) 是 \(R\) 的理想。则（加法）商群 \(R/I\) 在二元运算
\[
(r+I) + (s+I) = (r+s)+I, \qquad (r+I) \times (s+I) = (rs)+I
\]
（对所有 \(r, s \in R\)）下是环。反之，若 \(I\) 是任意使上述运算良定义的子群，则 \(I\) 是 \(R\) 的理想。
\end{proposition}

\begin{definition}
当 \(I\) 是 \(R\) 的理想时，带前一命题中运算的环 \(R/I\) 称为 \(R\) 对 \(I\) 的\textbf{商环}。
\end{definition}

\begin{theorem}\label{thm:7.7}
（1）（环的第一同构定理）若 \(\varphi : R \to S\) 是环的同态，则 \(\varphi\) 的核是 \(R\) 的理想，\(\varphi\) 的像是 \(S\) 的子环，且 \(R/\ker\varphi\) 作为环同构于 \(\varphi(R)\).

（2）若 \(I\) 是 \(R\) 的任意理想，则由
\[
\pi : R \to R/I, \quad r \mapsto r+I
\]
定义的映射是核为 \(I\) 的满环同态（这个同态称为 \(R\) 到 \(R/I\) 上的\textbf{自然投影}）。因此每个理想都是环同态的核，反之亦然。
\end{theorem}

\begin{proof}
这只是把前面的计算汇总起来。若 \(I\) 是 \(\varphi\) 的核，则 \(I\) 的（加法）陪集恰好是 \(\varphi\) 的纤维。特别地，陪集 \(r+I\)、\(s+I\) 和 \(rs+I\) 分别是 \(\varphi\) 在 \(\varphi(r)\)、\(\varphi(s)\) 和 \(\varphi(rs)\) 上方的纤维。由于 \(\varphi\) 是环同态，\(\varphi(r)\varphi(s) = \varphi(rs)\)，故 \((r+I)(s+I) = rs+I\). 陪集的乘法良定义，故 \(I\) 是理想且 \(R/I\) 是环。对应 \(r+I \mapsto \varphi(r)\) 是环 \(R/I\) 与 \(\varphi(R)\) 之间的保持加法与乘法的双射，从而是环同构。

若 \(I\) 是任意理想，则 \(R/I\) 是环（特别是阿贝尔群），映射 \(\pi : r \mapsto r+I\) 是核为 \(I\) 的群同态。还需验证 \(\pi\) 是环同态。这由 \(R/I\) 中乘法的定义立即得到：
\[
\pi(rs) = rs+I = (r+I)(s+I) = \pi(r)\pi(s).
\]
\end{proof}

与群一样，我们常常用横线记号表示模 \(I\) 的约化：\(\bar{r} = r+I\). 用这个记号，商环 \(R/I\) 中的加法与乘法简单地变成 \(\bar{r}+\bar{s} = \overline{r+s}\) 和 \(\bar{r}\,\bar{s} = \overline{rs}\).

\begin{example}
设 \(R\) 是环。

（1）子环 \(R\) 和 \(\{0\}\) 是理想。若 \(I \neq R\)，则称理想 \(I\) 是\textbf{真的}。理想 \(\{0\}\) 称为\textbf{平凡理想}，记作 0.

（2）立即看出对任意 \(n \in \mathbb{Z}\)，\(n\mathbb{Z}\) 是 \(\mathbb{Z}\) 的理想，且这些是 \(\mathbb{Z}\) 仅有的理想，因为特别地这些是 \(\mathbb{Z}\) 仅有的子群。相应的商环是 \(\mathbb{Z}/n\mathbb{Z}\)（这解释了记号的选择，且我们现在已经证明它是环），在第 0 章引入。例如若 \(n = 15\)，则 \(\mathbb{Z}/15\mathbb{Z}\) 的元素是陪集 \(\bar{0}, \bar{1}, \ldots, \overline{13}, \overline{14}\). 要在商中相加（或相乘），只需为两个陪集选取任意代表元，在整数 \(\mathbb{Z}\) 中把这些代表元相加（相乘），然后取包含这个和（积）的相应陪集。例如 \(\bar{7} + \overline{11} = \overline{18}\)，而 \(\overline{18} = \bar{3}\)，故在 \(\mathbb{Z}/15\mathbb{Z}\) 中 \(\bar{7}+\overline{11} = \bar{3}\). 类似地，在 \(\mathbb{Z}/15\mathbb{Z}\) 中 \(\bar{7}\cdot\overline{11} = \overline{77} = \bar{2}\). 我们也可以写成 \(7+11 \equiv 3 \pmod{15}\)、\(7 \cdot 11 \equiv 2 \pmod{15}\).

自然投影 \(\mathbb{Z} \to \mathbb{Z}/n\mathbb{Z}\) 称为\textbf{模 \(n\) 约化}，将在这些例子末尾进一步讨论。

（3）设 \(R = \mathbb{Z}[x]\) 是整系数 \(x\) 的多项式环。设 \(I\) 是所有项的次数至少为 2（即没有 0 次或 1 次项）的多项式连同零多项式的集合。则 \(I\) 是理想：两个这样的多项式之和的项次数仍至少为 2，而项次数至少为 2 的多项式与任意多项式之积的项次数仍至少为 2. 两个多项式 \(p(x), q(x)\) 在 \(I\) 的同一陪集中当且仅当它们相差一个项次数至少为 2 的多项式，即当且仅当 \(p(x)\) 和 \(q(x)\) 有相同的常数项和一次项。例如多项式 \(3+5x+x^3+x^5\) 和 \(3+5x-x^4\) 在 \(I\) 的同一陪集中。由此容易推出商 \(R/I\) 的完全代表元集由次数至多为 1 的多项式 \(a+bx\) 给出。

商中的加法与乘法同样用代表元进行。例如
\[
(1+3x) + (-4+5x) = -3+5\bar{x}
\]
且
\[
(1+3x)(-4+5x) = (-4-7x+15x^2) = -4-\overline{7}x.
\]
注意在这个商环 \(R/I\) 中例如 \(\bar{x}\cdot\bar{x} = \overline{x^2} = 0\)，故 \(R/I\) 有零因子，尽管 \(R = \mathbb{Z}[x]\) 没有。

（4）设 \(A\) 是环，\(X\) 是任意非空集合，\(R\) 是所有从 \(X\) 到 \(A\) 的函数构成的环。对每个固定的 \(c \in X\)，由
\[
E_c : R \to A, \quad E_c(f) = f(c)
\]
定义的映射（称为\textbf{在 \(c\) 处的求值}）是环同态，因为 \(R\) 中的运算是函数的逐点加法与乘法。\(E_c\) 的核由 \(\{f \in R \mid f(c) = 0\}\) 给出（即在 \(c\) 处为零的从 \(X\) 到 \(A\) 的函数集合）。此外 \(E_c\) 是满射：给定任意 \(a \in A\)，常函数 \(f(x) = a\) 在 \(c\) 处求值映到 \(a\). 故 \(R/\ker E_c \cong A\).

类似地，设 \(X\) 是 \(\mathbb{R}\) 中的闭区间 \([0,1]\)，\(R\) 是 \([0,1]\) 上所有连续实值函数构成的环。对每个 \(c \in [0,1]\)，在 \(c\) 处求值是满环同态（因为 \(R\) 含常函数），故 \(R/\ker E_c \cong \mathbb{R}\). \(E_c\) 的核是所有图像在 \(c\) 处穿过 \(x\) 轴的连续函数构成的理想。更一般地，\(E_c\) 在实数 \(y_0\) 上方的纤维是所有经过点 \((c, y_0)\) 的连续函数的集合。

（5）由 \(p(x) \mapsto p(0)\)（在 0 处求值）定义的从多项式环 \(R[x]\) 到 \(R\) 的映射是环同态，其核是所有常数项为零（即 \(p(0) = 0\)）的多项式集合。我们可以把这个同态与从 \(R\) 到另一个环 \(S\) 的任何同态复合，得到从 \(R[x]\) 到 \(S\) 的环同态。例如设 \(R = \mathbb{Z}\)，考虑由复合 \(p(x) \mapsto p(0) \mapsto p(0) \bmod 2 \in \mathbb{Z}/2\mathbb{Z}\) 定义的 \(\mathbb{Z}[x] \to \mathbb{Z}/2\mathbb{Z}\) 同态。这个复合映射的核由 \(\{p(x) \in \mathbb{Z}[x] \mid p(0) \in 2\mathbb{Z}\}\) 给出，即所有常数项为偶数的整系数多项式集合。这个同态的另一条纤维是常数项为奇数的多项式陪集，正如我们先前确定的。由于这个同态显然满射，商环是 \(\mathbb{Z}/2\mathbb{Z}\).

（6）固定某个 \(n \in \mathbb{Z}\)（\(n \geq 2\)），考虑非交换环 \(M_n(R)\). 若 \(J\) 是 \(R\) 的任意理想，则 \(M_n(J)\)（元素取自 \(J\) 的 \(n \times n\) 矩阵）是 \(M_n(R)\) 的双边理想。这个理想是满同态 \(M_n(R) \to M_n(R/J)\) 的核，后者把矩阵的每个元素模 \(J\) 约化，即把每个元素 \(a_{ij}\) 映到 \(\overline{a_{ij}}\)（这里横线表示过渡到 \(R/J\)）。例如当 \(n = 3\)、\(R = \mathbb{Z}\) 时，元素全为偶数的 \(3 \times 3\) 矩阵是 \(M_3(\mathbb{Z})\) 的双边理想 \(M_3(2\mathbb{Z})\)，且商 \(M_3(\mathbb{Z})/M_3(2\mathbb{Z})\) 同构于 \(M_3(\mathbb{Z}/2\mathbb{Z})\). 若环 \(R\) 有单位元，则下面的习题表明 \(M_n(R)\) 的每个双边理想都形如 \(M_n(J)\)，其中 \(J\) 是 \(R\) 的某个双边理想。

（7）设 \(R\) 是带 1 的交换环，\(G = \{g_1, \ldots, g_n\}\) 是有限群。由 \(\sum_{i=1}^{n} a_ig_i \mapsto \sum_{i=1}^{n} a_i\) 定义的从群环 \(RG\) 到 \(R\) 的映射容易看出是同态，称为\textbf{增广映射}。增广映射的核，即\textbf{增广理想}，是 \(RG\) 中系数之和为零的元素集合。例如 \(g_i - g_1\) 对所有 \(i, j\) 都是增广理想的元素。由于增广映射满射，商环同构于 \(R\).

\(RG\) 中的另一个理想是 \(\{\sum_{i=1}^{n} ag_i \mid a \in R\}\)，即系数全相等的形式和（等价地，元素 \(g_1 + \cdots + g_n\) 的所有 \(R\)-倍数）。

（8）设 \(R\) 是带单位元 \(1 \neq 0\) 的交换环，\(n \in \mathbb{Z}\)（\(n \geq 2\)）。我们展出 \(M_n(R)\) 中的一些单边理想。对每个 \(j \in \{1, 2, \ldots, n\}\)，设 \(L_j\) 是 \(M_n(R)\) 中第 \(j\) 列元素任意、其余列全为零的所有 \(n \times n\) 矩阵的集合。显然 \(L_j\) 对减法封闭。由矩阵乘法的定义直接可得：对任意矩阵 \(T \in M_n(R)\) 和任意 \(A \in L_j\)，乘积 \(TA\) 的第 \(i\) 列（\(i \neq j\)）元素为零。这表明 \(L_j\) 是 \(M_n(R)\) 的左理想。此外，\(L_j\) 不是右理想（因而不是双边理想）。为看出这一点，设 \(E_{pq}\) 是第 \(p\) 行第 \(q\) 列元素为 1、其余为零的矩阵（\(p, q \in \{1, \ldots, n\}\)）。则 \(E_{ij} \in L_j\) 但若 \(i \neq j\) 则 \(E_{ij}E_{ji} = E_{ii} \notin L_j\)，故 \(L_j\) 对右乘任意环元素不封闭。类似的论证表明：若 \(R_j\) 是 \(M_n(R)\) 中第 \(j\) 行元素任意、其余行全为零的所有 \(n \times n\) 矩阵的集合，则 \(R_j\) 是右理想但不是左理想。这些单边理想将在第 VI 部分起重要作用。
\end{example}

\begin{example}[约化同态]
从 \(\mathbb{Z}\) 到 \(\mathbb{Z}/n\mathbb{Z}\) 的、通过取出 \(\mathbb{Z}\) 的理想 \(n\mathbb{Z}\) 作商得到的标准投影映射通常称为“模 \(n\) 约化”。它是环同态这一事实对初等数论有重要推论。例如，假设我们试图求解方程
\[
x^2 + y^2 = 3z^2
\]
的整数解 \(x, y, z\)（这类问题常称为丢番图方程，以 Diophantus 命名，他是最早系统考察方程整数解存在性的人之一）。假设这样的整数存在。首先注意我们可以假设 \(x, y, z\) 没有公因子，否则我们可以用这个公因子的平方除这个方程，得到另一组比初始解更小的整数解。这个方程简单地陈述了环 \(\mathbb{Z}\) 中这些元素之间的一个关系。因此同样的关系在任何商环中也必成立。特别地，这个关系在 \(\mathbb{Z}/n\mathbb{Z}\) 中对任意整数 \(n\) 都成立。选择 \(n = 4\) 特别有效，理由如下：模 4 的平方只是 \(0^2, 1^2, 2^2, 3^2\)，即 \(0, 1 \pmod 4\). 把这个方程模 4 读取（即在商环 \(\mathbb{Z}/4\mathbb{Z}\) 中考虑这个方程），我们必须有
\[
\left\{\begin{smallmatrix}0\\1\end{smallmatrix}\right\} + \left\{\begin{smallmatrix}0\\1\end{smallmatrix}\right\} = 3\left\{\begin{smallmatrix}0\\1\end{smallmatrix}\right\} \pmod 4,
\]
其中例如 \(\left\{\begin{smallmatrix}0\\1\end{smallmatrix}\right\}\) 表示可以取 0 或 1. 检查少数几种可能性表明每次都必须取 0. 这意味着 \(x, y, z\) 每一个都必是偶数（奇数的平方给我们 \(1 \bmod 4\)）。但这与这些整数无公因子的假设矛盾，从而表明这个方程没有非零整数解。

注意即使解存在，这个技巧也给出关于解的可能剩余类的信息（因为我们同样可以考察模 \(n\) 而非模 4 的可能性），且注意对每个 \(n\) 的选择我们只需解一个有限问题，因为模 \(n\) 的剩余类只有有限多个。连同中国剩余定理（在 7.6 节描述），我们随后可以确定模非常大的整数时的可能解，这大大有助于数值求解（当解存在时）。我们还注意到这个技巧有许多局限——例如存在模每个整数都有解但没有整数解的方程。一个容易写出的例子是方程
\[
(x^2 - 13)(x^2 - 17)(x^2 - 221) = 0.
\]

作为这个技巧的最后一个例子，我们提到对素数 \(p\)，由把系数模 \(p\) 约化给出的从整系数多项式环 \(\mathbb{Z}[x]\) 到系数在 \(\mathbb{Z}/p\mathbb{Z}\) 中的多项式环 \(\mathbb{Z}/p\mathbb{Z}[x]\) 的映射是环同态。这个约化例子将在第 9 章用于判定多项式能否分解。
\end{example}

下面的定理给出环的其余同构定理。其中每一个都可以如下证明：首先用群论中相应的定理得到加法群的同构（或第四同构定理情形下群的对应），然后验证这个群同构（或对应）是乘性映射，从而定义环同构。在每种情形验证都可由商环中乘法的定义立即得到。例如，给出下面（2）中同构的映射由 \(\varphi : r+I \mapsto r+J\) 定义。这个映射是乘性的，因为由商环 \(R/I\) 中乘法的定义 \((r_1+I)(r_2+I) = r_1r_2+I\)，而 \(r_1r_2+I \mapsto r_1r_2+J = (r_1+J)(r_2+J)\)（由商环 \(R/J\) 中乘法的定义），即 \(\varphi(r_1r_2) = \varphi(r_1)\varphi(r_2)\). 定理其余部分的证明类似。

\begin{theorem}\label{thm:7.8}
设 \(R\) 是环。

（1）（环的第二同构定理）设 \(A\) 是子环，\(B\) 是 \(R\) 的理想。则 \(A+B = \{a+b \mid a \in A,\ b \in B\}\) 是 \(R\) 的子环，\(A \cap B\) 是 \(A\) 的理想，且 \((A+B)/B \cong A/(A \cap B)\).

（2）（环的第三同构定理）设 \(I\) 和 \(J\) 是 \(R\) 的理想且 \(I \subseteq J\). 则 \(J/I\) 是 \(R/I\) 的理想，且 \((R/I)/(J/I) \cong R/J\).

（3）（环的第四或格子同构定理）设 \(I\) 是 \(R\) 的理想。对应 \(A \leftrightarrow A/I\) 是包含 \(I\) 的 \(R\) 的子环 \(A\) 的集合与 \(R/I\) 的子环的集合之间保持包含关系的双射。此外，\(A\)（含 \(I\) 的子环）是 \(R\) 的理想当且仅当 \(A/I\) 是 \(R/I\) 的理想。
\end{theorem}

设 \(R = \mathbb{Z}\)，\(I\) 是理想 \(12\mathbb{Z}\). 商环 \(\bar{R} = R/I = \mathbb{Z}/12\mathbb{Z}\) 有理想 \(\bar{R}\)、\(2\mathbb{Z}/12\mathbb{Z}\)、\(3\mathbb{Z}/12\mathbb{Z}\)、\(4\mathbb{Z}/12\mathbb{Z}\)、\(6\mathbb{Z}/12\mathbb{Z}\) 和 \(0 = 12\mathbb{Z}/12\mathbb{Z}\)，分别对应于含 \(I\) 的 \(R = \mathbb{Z}\)、\(2\mathbb{Z}\)、\(3\mathbb{Z}\)、\(4\mathbb{Z}\)、\(6\mathbb{Z}\)、\(12\mathbb{Z} = I\) 各理想。

若 \(I\) 和 \(J\) 是环 \(R\) 中的理想，则 \(a \in I\)、\(b \in J\) 的和 \(a+b\) 的集合不仅是 \(R\) 的子环（如环的第二同构定理），而且是 \(R\) 的理想（这个集合显然对加法封闭，且 \(r(a+b) = ra+rb \in I+J\)，因为 \(ra \in I\)、\(rb \in J\)）。我们还可以定义两个理想的积：

\begin{definition}
设 \(I\) 和 \(J\) 是 \(R\) 的理想。

（1）定义 \(I\) 与 \(J\) 的\textbf{和}为 \(I+J = \{a+b \mid a \in I,\ b \in J\}\).

（2）定义 \(I\) 与 \(J\) 的\textbf{积}（记作 \(IJ\)）为形如 \(ab\)（\(a \in I\)，\(b \in J\)）的元素的全部有限和构成的集合。

（3）对任意 \(n \geq 1\)，定义 \(I\) 的 \(n\) 次幂 \(I^n\) 为形如 \(a_1a_2\cdots a_n\)（\(a_i \in I\)）的元素的全部有限和构成的集合。等价地，\(I^n\) 通过定义 \(I^1 = I\)、\(I^n = II^{n-1}\)（\(n = 2, 3, \ldots\)）归纳定义。
\end{definition}

容易看出理想 \(I\) 与 \(J\) 的和 \(I+J\) 是 \(R\) 的包含 \(I\) 和 \(J\) 两者的最小理想，积 \(IJ\) 是含于 \(I \cap J\) 中的理想（但可能严格更小，参见习题）。还要注意积理想 \(IJ\) 的元素是来自 \(I\) 和 \(J\) 的元素之积 \(ab\) 的有限和。仅由 \(I\) 和 \(J\) 中元素之积组成的集合 \(\{ab \mid a \in I,\ b \in J\}\) 一般不对加法封闭，故一般不是理想。

\begin{example}
（1）设 \(I = 6\mathbb{Z}\)、\(J = 10\mathbb{Z}\) 是 \(\mathbb{Z}\) 中的理想。则 \(I+J\) 由所有形如 \(6x+10y\)（\(x, y \in \mathbb{Z}\)）的整数组成。由于每个这样的整数都被 2 整除，理想 \(I+J\) 包含在 \(2\mathbb{Z}\) 中。另一方面，\(2 = 6(2) + 10(-1)\) 表明理想 \(I+J\) 包含理想 \(2\mathbb{Z}\)，故 \(6\mathbb{Z}+10\mathbb{Z} = 2\mathbb{Z}\). 一般地，\(m\mathbb{Z}+n\mathbb{Z} = d\mathbb{Z}\)，其中 \(d\) 是 \(m\) 与 \(n\) 的最大公因子。积 \(IJ\) 由形如 \((6x)(10y)\)（\(x, y \in \mathbb{Z}\)）的元素的全部有限和组成，它显然给出理想 \(60\mathbb{Z}\).

（2）设 \(I\) 是 \(\mathbb{Z}[x]\) 中常数项为偶数的整系数多项式构成的理想（参见例 5）。两个多项式 2 和 \(x\) 含于 \(I\) 中，故 \(4 = 2 \cdot 2\) 和 \(x^2 = x \cdot x\) 都是积理想 \(I^2 = II\) 的元素，它们的和 \(x^2+4\) 也是。然而容易验证 \(x^2+4\) 不能写成 \(I\) 的两个元素的单个乘积 \(p(x)q(x)\).
\end{example}

\subsection*{习题}
设 \(R\) 是带单位元 \(1 \neq 0\) 的环。

\begin{enumerate}
\item 证明环 \(2\mathbb{Z}\) 与 \(3\mathbb{Z}\) 不同构。

\item 证明环 \(\mathbb{Z}[x]\) 与 \(\mathbb{Q}[x]\) 不同构。

\item 求 \(\mathbb{Z}\) 的所有同态像。

\item 求从 \(\mathbb{Z}\) 到 \(\mathbb{Z}/30\mathbb{Z}\) 的所有环同态。每种情形描述核与像。

\item 描述从环 \(\mathbb{Z} \times \mathbb{Z}\) 到 \(\mathbb{Z}\) 的所有环同态。每种情形描述核与像。

\item 判断下列哪些是从 \(M_2(\mathbb{Z})\) 到 \(\mathbb{Z}\) 的环同态：

（a）\(\begin{pmatrix} a & b \\ c & d \end{pmatrix} \mapsto a\)（投影到 (1,1) 元素）；

（b）\(\begin{pmatrix} a & b \\ c & d \end{pmatrix} \mapsto a+d\)（矩阵的迹）；

（c）\(\begin{pmatrix} a & b \\ c & d \end{pmatrix} \mapsto ad-bc\)（矩阵的行列式）。

\item 设 \(R = \left\{\begin{pmatrix} a & b \\ 0 & d \end{pmatrix} \ \middle|\ a, b, d \in \mathbb{Z}\right\}\) 是 \(M_2(\mathbb{Z})\) 的上三角矩阵子环。证明映射
\[
\varphi : R \to \mathbb{Z} \times \mathbb{Z}, \quad \varphi : \begin{pmatrix} a & b \\ 0 & d \end{pmatrix} \mapsto (a, d)
\]
是满同态，并描述它的核。

\item 判断下列哪些是环 \(\mathbb{Z} \times \mathbb{Z}\) 的理想：

（a）\(\{(a, a) \mid a \in \mathbb{Z}\}\)；（b）\(\{(2a, 2b) \mid a, b \in \mathbb{Z}\}\)；（c）\(\{(2a, 0) \mid a \in \mathbb{Z}\}\)；（d）\(\{(a, -a) \mid a \in \mathbb{Z}\}\).

\item 判断第 1 节习题 6 中的哪些集合是从 \([0,1]\) 到 \(\mathbb{R}\) 的所有函数环的理想。

\item 判断下列哪些是环 \(\mathbb{Z}[x]\) 的理想：

（a）所有常数项是 3 的倍数的多项式集合；

（b）所有 \(x^2\) 的系数是 3 的倍数的多项式集合；

（c）所有常数项、\(x\) 的系数和 \(x^2\) 的系数都为零的多项式集合；

（d）\(\mathbb{Z}[x^2]\)（即只出现 \(x\) 的偶次幂的多项式）；

（e）系数之和为零的多项式集合；

（f）满足 \(p'(0) = 0\) 的多项式 \(p(x)\) 的集合，其中 \(p'(x)\) 是 \(p(x)\) 对 \(x\) 的通常一阶导数。

\item 设 \(R\) 是闭区间 \([0,1]\) 上所有连续实值函数构成的环。证明由 \(\varphi(f) = \int_0^1 f(t)\,dt\) 定义的映射 \(\varphi : R \to \mathbb{R}\) 是加法群的同态但不是环同态。

\item 设 \(D\) 是 \(\mathbb{Z}\) 中不是完全平方的整数，\(S = \left\{\begin{pmatrix} a & b \\ Db & a \end{pmatrix} \ \middle|\ a, b \in \mathbb{Z}\right\}\).

（a）证明 \(S\) 是 \(M_2(\mathbb{Z})\) 的子环。

（b）若 \(D\) 在 \(\mathbb{Z}\) 中不是完全平方，证明由 \(\varphi(a + b\sqrt{D}) = \begin{pmatrix} a & b \\ Db & a \end{pmatrix}\) 定义的映射 \(\varphi : \mathbb{Z}[\sqrt{D}] \to S\) 是环同构。

（c）若 \(D \equiv 1 \pmod 4\) 无平方因子，证明集合 \(\left\{\begin{pmatrix} a & b \\ \frac{(D-1)b}{4} & a+b \end{pmatrix} \ \middle|\ a, b \in \mathbb{Z}\right\}\) 是 \(M_2(\mathbb{Z})\) 的子环，且同构于二次整数环 \(\mathcal{O}\).

\item 证明环 \(M_2(\mathbb{R})\) 含有一个同构于 \(\mathbb{C}\) 的子环。

\item 证明环 \(M_4(\mathbb{R})\) 含有一个同构于实 Hamilton 四元数 \(\mathbb{H}\) 的子环。

\item 设 \(X\) 是非空集合，\(\mathcal{P}(X)\) 是第 1 节习题 21 中定义的 \(X\) 的所有子集的布尔环。设 \(R\) 是所有从 \(X\) 到 \(\mathbb{Z}/2\mathbb{Z}\) 的函数构成的环。对每个 \(A \in \mathcal{P}(X)\) 定义函数
\[
\chi_A : X \to \mathbb{Z}/2\mathbb{Z}, \quad \chi_A(x) = \begin{cases} 1, & \text{若 } x \in A, \\ 0, & \text{若 } x \notin A \end{cases}
\]
（\(\chi_A\) 称为 \(A\) 的、取值在 \(\mathbb{Z}/2\mathbb{Z}\) 中的\textbf{特征函数}）。证明由 \(A \mapsto \chi_A\) 定义的映射 \(\mathcal{P}(X) \to R\) 是环同构。

\item 设 \(\varphi : R \to S\) 是满环同态。证明 \(R\) 的中心的像包含在 \(S\) 的中心中（参见第 1 节习题 7）。

\item 设 \(R\) 和 \(S\) 是带单位元的非零环，其单位元分别记作 \(1_R\) 和 \(1_S\). 设 \(\varphi : R \to S\) 是非零环同态。

（a）证明若 \(\varphi(1_R) \neq 1_S\)，则 \(\varphi(1_R)\) 是 \(S\) 中的零因子。由此推出：若 \(S\) 是整环，则从 \(R\) 到 \(S\) 的每个环同态都把 \(R\) 的单位元送到 \(S\) 的单位元。

（b）证明若 \(\varphi(1_R) = 1_S\)，则对 \(R\) 的每个单位 \(u\)，\(\varphi(u)\) 是 \(S\) 中的单位且 \(\varphi(u^{-1}) = \varphi(u)^{-1}\).

\item （a）若 \(I\) 和 \(J\) 是 \(R\) 的理想，证明它们的交 \(I \cap J\) 也是 \(R\) 的理想。

（b）证明任意非空族理想的交也是理想（不要假设该族可数）。

\item 证明：若 \(I_1 \subseteq I_2 \subseteq \cdots\) 是 \(R\) 的理想，则 \(\bigcup_{n} I_n\) 是 \(R\) 的理想。

\item 设 \(I\) 是 \(R\) 的理想，\(S\) 是 \(R\) 的子环。证明 \(I \cap S\) 是 \(S\) 的理想。举例说明环 \(R\) 的子环 \(S\) 的理想不必都形如 \(I \cap S\)，其中 \(I\) 是 \(R\) 的某个理想。

\item 证明 \(M_n(R)\) 的每个（双边）理想都等于 \(M_n(J)\)，其中 \(J\) 是 \(R\) 的某个（双边）理想。〔用第 2 节习题 6(c) 首先证明 \(M_n(R)\) 的理想的矩阵元素集合构成 \(R\) 的理想。〕

\item 设 \(a\) 是环 \(R\) 的元素。

（a）证明 \(\{x \in R \mid ax = 0\}\) 是右理想，\(\{y \in R \mid ya = 0\}\) 是左理想（分别称为 \(a\) 在 \(R\) 中的右、左\textbf{零化子}）。

（b）证明若 \(L\) 是 \(R\) 的左理想，则 \(\{x \in R \mid xa = 0 \text{ 对所有 } a \in L\}\) 是双边理想（称为 \(L\) 在 \(R\) 中的左零化子）。

\item 设 \(S\) 是 \(R\) 的子环，\(I\) 是 \(R\) 的理想。证明若 \(S \cap I = 0\)，则 \(\bar{S} \cong S\)，其中横线表示过渡到 \(R/I\).

\item 设 \(\varphi : R \to S\) 是环同态。

（a）证明若 \(J\) 是 \(S\) 的理想，则 \(\varphi^{-1}(J)\) 是 \(R\) 的理想。把它应用于 \(R\) 是 \(S\) 的子环、\(\varphi\) 是包含同态的特殊情形，推出若 \(J\) 是 \(S\) 的理想则 \(J \cap R\) 是 \(R\) 的理想。

（b）证明若 \(\varphi\) 满射且 \(I\) 是 \(R\) 的理想，则 \(\varphi(I)\) 是 \(S\) 的理想。举一个 \(\varphi\) 不满射时这一结论不成立的例子。

\item 假设 \(R\) 是带 1 的交换环。证明二项式定理
\[
(a+b)^n = \sum_{k=0}^{n}\binom{n}{k}a^kb^{n-k}
\]
在 \(R\) 中成立，其中二项式系数 \(\binom{n}{k}\) 在 \(R\) 中解释为 \(R\) 的单位元 1 的 \(\binom{n}{k}\) 次和 \(1+1+\cdots+1\).

\item 环 \(R\) 的\textbf{特征}是使 \(1+1+\cdots+1 = 0\)（\(n\) 个 1）在 \(R\) 中成立的最小正整数 \(n\)；若这样的整数不存在，则称 \(R\) 的特征为 0. 例如 \(\mathbb{Z}/n\mathbb{Z}\) 对每个正整数 \(n\) 是特征为 \(n\) 的环，\(\mathbb{Z}\) 是特征为 0 的环。

（a）证明由
\[
k \mapsto \begin{cases} 1+1+\cdots+1\ (k \text{ 个 } 1), & \text{若 } k > 0, \\ 0, & \text{若 } k = 0, \\ -1-1-\cdots-1\ (-k \text{ 个}), & \text{若 } k < 0 \end{cases}
\]
定义的映射 \(\mathbb{Z} \to R\) 是环同态，其核是 \(n\mathbb{Z}\)，其中 \(n\) 是 \(R\) 的特征（这解释了“特征 0”这一术语的用法，而不是对没有 1 的和为零的环使用古语“特征 \(\infty\)”）。

（b）确定环 \(\mathbb{Q}\)、\(\mathbb{Z}[x]\)、\(\mathbb{Z}/n\mathbb{Z}[x]\) 的特征。

（c）证明若 \(p\) 是素数且 \(R\) 是特征为 \(p\) 的交换环，则对所有 \(a, b \in R\) 有 \((a+b)^p = a^p+b^p\).

\item 证明非零布尔环有特征 2（参见第 1 节习题 15）。

\item 证明整环的特征为 \(p\)，其中 \(p\) 是素数或 0（参见习题 26）。

\item 设 \(R\) 是交换环。回忆（参见第 1 节习题 13）元素 \(x \in R\) 是幂零的，若对某个 \(n \in \mathbb{Z}^+\) 有 \(x^n = 0\). 证明幂零元素的集合构成一个理想——称为 \(R\) 的\textbf{幂零根}，记作 \(\mathfrak{N}(R)\).〔用二项式定理证明 \(\mathfrak{N}(R)\) 对加法封闭。〕

\item 证明：若 \(R\) 是交换环且 \(\mathfrak{N}(R)\) 是其幂零根（参见前一习题），则零是 \(R/\mathfrak{N}(R)\) 中唯一的幂零元素，即证明 \(\mathfrak{N}(R/\mathfrak{N}(R)) = 0\).

\item 证明元素 \(\begin{pmatrix} 0 & 1 \\ 0 & 0 \end{pmatrix}\) 和 \(\begin{pmatrix} 0 & 0 \\ 1 & 0 \end{pmatrix}\) 是 \(M_2(\mathbb{Z})\) 的幂零元素，它们的和不是幂零的（注意这两个矩阵不交换）。由此推出非交换环 \(M_2(\mathbb{Z})\) 中幂零元素的集合不是理想。

\item 设 \(\varphi : R \to S\) 是环的同态。证明若 \(x\) 是 \(R\) 的幂零元素，则 \(\varphi(x)\) 在 \(S\) 中是幂零的。

\item 假设 \(R\) 交换。设 \(p(x) = a_nx^n + a_{n-1}x^{n-1} + \cdots + a_1x + a_0\) 是多项式环 \(R[x]\) 的元素。

（a）证明 \(p(x)\) 是 \(R[x]\) 中的单位当且仅当 \(a_0\) 是单位且 \(a_1, a_2, \ldots, a_n\) 在 \(R\) 中幂零。〔见第 1 节习题 14。〕

（b）证明 \(p(x)\) 在 \(R[x]\) 中幂零当且仅当 \(a_0, a_1, \ldots, a_n\) 都是 \(R\) 的幂零元素。

\item 设 \(I, J\) 是 \(R\) 的理想。

（a）证明 \(I+J\) 是 \(R\) 的包含 \(I\) 和 \(J\) 两者的最小理想。

（b）证明 \(IJ\) 是含于 \(I \cap J\) 中的理想。

（c）举一个 \(IJ \neq I \cap J\) 的例子。

（d）证明若 \(R\) 交换且 \(I+J = R\)，则 \(IJ = I \cap J\).

\item 设 \(I, J, K\) 是 \(R\) 的理想。

（a）证明 \(I(J+K) = IJ+IK\) 且 \((I+J)K = IK+JK\).

（b）证明若 \(J \subseteq I\) 则 \(I \cap (J+K) = J + (I \cap K)\).

\item 证明：若 \(I\) 是 \(\mathbb{Z}[x]\) 中常数项为零的所有多项式构成的理想，则
\[
I^n = \{a_nx^n + a_{n+1}x^{n+1} + \cdots + a_{n+m}x^{n+m} \mid a_i \in \mathbb{Z},\ m \geq 0\}
\]
是第一个非零项次数至少为 \(n\) 的多项式集合。

\item 理想 \(N\) 称为\textbf{幂零的}，若对某个 \(n \geq 1\)，\(N^n\) 是零理想。证明理想 \(p\mathbb{Z}/p^n\mathbb{Z}\) 是环 \(\mathbb{Z}/p^n\mathbb{Z}\) 中的幂零理想。
\end{enumerate}

\section{理想的性质}

本节自始至终 \(R\) 是带单位元 \(1 \neq 0\) 的环。

\begin{definition}
设 \(A\) 是环 \(R\) 的任意子集。

（1）用 \((A)\) 表示 \(R\) 的包含 \(A\) 的最小理想，称为 \(A\) \textbf{生成的理想}。

（2）用 \(RA\) 表示所有形如 \(ra\)（\(r \in R\)，\(a \in A\)）的元素的有限和构成的集合，即
\[
RA = \{r_1a_1 + r_2a_2 + \cdots + r_na_n \mid r_i \in R,\ a_i \in A,\ n \in \mathbb{Z}^+\}
\]
（约定若 \(A = \varnothing\) 则 \(RA = 0\)）。类似地，\(AR = \{a_1r_1 + a_2r_2 + \cdots + a_nr_n \mid r_i \in R,\ a_i \in A,\ n \in \mathbb{Z}^+\}\)，\(RAR = \{r_1a_1r_1' + r_2a_2r_2' + \cdots + r_na_nr_n' \mid r_i, r_i' \in R,\ a_i \in A,\ n \in \mathbb{Z}^+\}\).

（3）由单个元素生成的理想称为\textbf{主理想}。

（4）由有限集生成的理想称为\textbf{有限生成理想}。
\end{definition}

当 \(A = \{a\}\) 或 \(\{a_1, a_2, \ldots\}\) 等时，我们去掉了集合括号，直接写 \((a)\)、\((a_1, a_2, \ldots)\) 表示 \((A)\).

由环的子集生成的理想的概念类似于由群的子集生成的子群（2.4 节）。由于 \(R\) 的任意非空族理想的交也是理想（参见第 3 节习题 18），且 \(A\) 总含于至少一个理想（即 \(R\)），我们有
\[
(A) = \bigcap_{\substack{I \text{ 是理想} \\ A \subseteq I}} I,
\]
即 \((A)\) 是所有包含集合 \(A\) 的 \(R\) 的理想的交。

由 \(A\) 生成的左理想是所有包含 \(A\) 的 \(R\) 的左理想的交。这个左理想由把 \(A\) 对定义左理想的所有运算封闭得到。由定义立即看出 \(RA\) 对加法和左乘任意环元素封闭。由于 \(R\) 有单位元，\(RA\) 包含 \(A\). 故 \(RA\) 是包含 \(A\) 的 \(R\) 的左理想。反之，任何包含 \(A\) 的左理想必包含所有形如 \(ra\)（\(r \in R\)，\(a \in A\)）的元素的有限和，从而必包含 \(RA\). 因此 \(RA\) 恰好是由 \(A\) 生成的左理想。类似地，\(AR\) 是由 \(A\) 生成的右理想，\(RAR\) 是由 \(A\) 生成的（双边）理想。特别地，
\begin{center}
若 \(R\) 交换，则 \(RA = AR = RAR = (A)\).
\end{center}

当 \(R\) 是交换环且 \(a \in R\) 时，由 \(a\) 生成的主理想 \((a)\) 就是 \(a\) 的所有 \(R\)-倍数的集合。然而若 \(R\) 不交换，集合 \(\{ras \mid r, s \in R\}\) 未必是由 \(a\) 生成的双边理想，因为它不必对加法封闭（此时由 \(a\) 生成的理想是理想 \(RaR\)，它由形如 \(ras\)（\(r, s \in R\)）的元素的全部有限和组成）。

在交换环中构造主理想是产生理想的一种特别简单的方式，类似于生成群的循环子群。注意元素 \(b \in R\) 属于理想 \((a)\) 当且仅当对某个 \(r \in R\) 有 \(b = ra\)，即当且仅当 \(b\) 是 \(a\) 的倍数，或者换句话说，\(a\) 在 \(R\) 中整除 \(b\). 又 \(b \in (a)\) 当且仅当 \((b) \subseteq (a)\). 因此理想之间（特别是主理想之间）的包含关系被认为捕捉了一般交换环的一部分算术。所有理想都是主理想的交换环是最容易研究的，它们将在第 8、9 章起重要作用。

\begin{example}
（1）平凡理想 0 和理想 \(R\) 都是主的：\(0 = (0)\)，\(R = (1)\).

（2）在 \(\mathbb{Z}\) 中，对所有整数 \(n\) 有 \(n\mathbb{Z} = \mathbb{Z}n = (n) = (-n)\). 因此我们对 \(n\mathbb{Z}\) 的记号与我们一直使用的 \(n\mathbb{Z}\) 的定义一致。如前一节所述，这些是 \(\mathbb{Z}\) 仅有的理想，故 \(\mathbb{Z}\) 的每个理想都是主的。对正整数 \(n, m\)，\(n\mathbb{Z} \supseteq m\mathbb{Z}\) 当且仅当 \(m\) 在 \(\mathbb{Z}\) 中整除 \(n\)，故包含 \(n\mathbb{Z}\) 的理想的格与 \(n\) 的因子的格相同。此外，两个非零整数 \(n, m\) 生成的理想是它们的最大公因子 \(d\) 生成的主理想：\((n, m) = (d)\). 因此 \((n, m)\) 作为 \(n\) 与 \(m\) 的最大公因子的记号与 \(n\) 和 \(m\) 生成的理想的同一记号一致（尽管 \(n\) 和 \(m\) 生成的理想的主生成元只确定到 \(\pm\) 号——我们可以通过选取非负生成元使它唯一）。特别地，\(n\) 与 \(m\) 互素当且仅当 \((n, m) = (1)\).

（3）我们证明 \(\mathbb{Z}[x]\) 中由 2 和 \(x\) 生成的理想 \((2, x)\) 不是主理想。注意 \((2, x) = \{2p(x) + xq(x) \mid p(x), q(x) \in \mathbb{Z}[x]\}\)，故这个理想恰好由常数项为偶数的整系数多项式组成（如前节例 5 所讨论）——特别地这是真理想。反设对某个 \(a(x) \in \mathbb{Z}[x]\) 有 \((2, x) = (a(x))\). 由于 \(2 \in (a(x))\)，必存在 \(p(x)\) 使 \(2 = p(x)a(x)\). \(p(x)a(x)\) 的次数等于 \(\deg p(x) + \deg a(x)\)，故 \(p(x)\) 和 \(a(x)\) 都必是常多项式，即整数。由于 2 是素数，\(a(x), p(x) \in \{\pm 1, \pm 2\}\). 若 \(a(x) = \pm 1\)，则每个多项式都是 \(a(x)\) 的倍数，与 \((a(x))\) 是真理想矛盾。唯一可能是 \(a(x) = \pm 2\). 但此时由 \(x \in (a(x)) = (2)\)（或 \((-2)\)）可知对某个整系数多项式 \(q(x)\) 有 \(x = 2q(x)\)，显然不可能。这个矛盾证明 \((2, x)\) 不是主的。

注意若不指明环，记号 \((A)\) 是含糊的：\(\mathbb{Q}[x]\) 中由 2 和 \(x\) 生成的理想是整个环 \((1)\)，因为它含元素 \(\frac{1}{2} \cdot 2 = 1\).

我们将在第 9 章看到对任意域 \(F\)，\(F[x]\) 的所有理想都是主的。

（4）若 \(R\) 是所有从闭区间 \([0,1]\) 到 \(\mathbb{R}\) 的函数构成的环，设 \(M\) 是理想 \(\{f \mid f(\tfrac12) = 0\}\)（在 \(\frac{1}{2}\) 处求值的核）。设 \(g(x)\) 是在 \(x = \frac{1}{2}\) 处为零、在其他所有点处为 1 的函数。则对所有 \(f \in M\) 有 \(f = fg\)，故 \(M\) 是以 \(g\) 为生成元的主理想。事实上，任何在 \(\frac{1}{2}\) 处为零且在其他所有点处非零的函数都是同一理想 \(M\) 的另一个生成元。

另一方面，若 \(R\) 是所有从 \([0,1]\) 到 \(\mathbb{R}\) 的连续函数构成的环，则 \(\{f \mid f(\tfrac12) = 0\}\) 不是主的，甚至不是有限生成的（参见习题）。

（5）若 \(G\) 是有限群、\(R\) 是带 1 的交换环，则增广理想 \(I\) 由集合 \(\{g - 1 \mid g \in G\}\) 生成，尽管这不必是极小的生成元集。例如若 \(G = \langle a \rangle\) 是循环群，则增广理想是以 \(a-1\) 为生成元的主理想。
\end{example}

\begin{proposition}\label{pro:7.9}
设 \(I\) 是 \(R\) 的理想。

\begin{enumerate}[label=(\arabic*)]
\item \(I = R\) 当且仅当 \(I\) 含有单位元（可逆元）。
\item 假设 \(R\) 交换。则 \(R\) 是域当且仅当它仅有的理想是 0 和 \(R\).
\end{enumerate}
\end{proposition}

\begin{proof}
（1）若 \(I = R\)，则 \(I\) 含单位元 1. 反之，若 \(u\) 是 \(I\) 中的单位、逆元为 \(v\)，则对任意 \(r \in R\) 有
\[
r = r\cdot 1 = r(vu) = (rv)u \in I,
\]
故 \(R = I\).

（2）环 \(R\) 是域当且仅当每个非零元素都是单位。若 \(R\) 是域，则每个非零理想都含一个单位，故由第一部分 \(R\) 是唯一的非零理想。反之，若 0 和 \(R\) 是 \(R\) 仅有的理想，设 \(u\) 是 \(R\) 的任意非零元素。由假设 \((u) = R\)，故 \(1 \in (u)\). 于是存在 \(v \in R\) 使 \(1 = vu\)，即 \(u\) 是单位。因此 \(R\) 的每个非零元素都是单位，故 \(R\) 是域。
\end{proof}

\begin{corollary}\label{cor:7.10}
若 \(R\) 是域，则从 \(R\) 到另一个环的任何非零环同态都是单射。
\end{corollary}

\begin{proof}
环同态的核是理想。非零同态的核是真理想，故由前一命题它是 0.
\end{proof}

这些结果表明域的理想结构是平凡的。我们通过同态研究代数结构的方法在域论（第 IV 部分）中仍将起基础作用，那时我们研究一个域到另一个域的单同态（嵌入）以及域的自同构（域到自身的同构）。

若 \(D\) 是带单位元 \(1 \neq 0\) 的环，且其仅有的左理想和仅有的右理想都是 0 和 \(D\)，则 \(D\) 是除环。反之，除环 \(D\) 中仅有的（左、右或双边）理想是 0 和 \(D\)，这给出 \(R\) 不交换时命题 \ref{pro:7.9}(2) 的一个类比（见习题）。然而若 \(F\) 是域，则对任意 \(n \geq 2\)，矩阵环 \(M_n(F)\) 中仅有的双边理想是 0 和 \(M_n(F)\)，尽管它不是除环（它确实有真非平凡左理想和右理想：参见 7.3 节），这表明命题 \ref{pro:7.9}(2) 对非交换环不成立。仅有的双边理想是 0 和整个环的环（称为\textbf{单环}）将在第 18 章研究。

一类重要的理想是那些不包含在任何其他真理想中的理想：

\begin{definition}
任意环 \(S\) 中的理想 \(M\) 称为\textbf{极大理想}，若 \(M \neq S\) 且包含 \(M\) 的理想只有 \(M\) 和 \(S\).
\end{definition}

一般环不一定有极大理想。例如取任意没有极大子群的阿贝尔群（例如 \(\mathbb{Q}\)——参见 4.3 节习题 16 和 6.1 节），通过定义对所有 \(a, b\) 有 \(ab = 0\) 把它做成平凡环。在这样的环中理想就是子群，故没有极大理想。零环没有极大理想，因此任何涉及极大理想的结果都迫使环非零。下一个命题表明带单位元 \(1 \neq 0\) 的环总有极大理想。像许多这样的存在性定理（例如有限生成群有极大子群、每个向量空间有基）一样，证明依赖于佐恩引理（参见附录 I）。然而在许多具体环中，极大理想的存在往往显然，与佐恩引理无关。

\begin{proposition}\label{pro:7.11}
带单位元的环中每个真理想都包含在一个极大理想中。
\end{proposition}

\begin{proof}
设 \(R\) 是带单位元的环，\(I\) 是真理想（故 \(R\) 不能是零环，即 \(1 \neq 0\)）。设 \(\mathcal{S}\) 是 \(R\) 的所有包含 \(I\) 的真理想的集合。则 \(\mathcal{S}\) 非空（\(I \in \mathcal{S}\)）且按包含关系偏序。若 \(\mathcal{C}\) 是 \(\mathcal{S}\) 中的链，定义 \(J\) 为 \(\mathcal{C}\) 中所有理想的并：
\[
J = \bigcup_{A \in \mathcal{C}} A.
\]
我们首先证明 \(J\) 是理想。当然 \(J\) 非空，因为 \(\mathcal{C}\) 非空——特别地 \(0 \in J\)，因为 0 在每个理想 \(A\) 中。若 \(a, b \in J\)，则存在 \(A, B \in \mathcal{C}\) 使 \(a \in A\)、\(b \in B\). 由链的定义，或者 \(A \subseteq B\) 或者 \(B \subseteq A\). 两种情形都有 \(a - b \in J\)，故 \(J\) 对减法封闭。由于每个 \(A \in \mathcal{C}\) 对左乘和右乘 \(R\) 中元素封闭，\(J\) 也如此。这证明 \(J\) 是理想。

若 \(J\) 不是真理想，则 \(1 \in J\). 此时由 \(J\) 的定义必有 \(1 \in A\) 对某个 \(A \in \mathcal{C}\). 这是矛盾，因为每个 \(A\) 是真理想（\(A \in \mathcal{C} \subseteq \mathcal{S}\)）。这证明每个链在 \(\mathcal{S}\) 中有上界。由佐恩引理，\(\mathcal{S}\) 有极大元，它因而就是包含 \(I\) 的极大（真）理想。
\end{proof}

对交换环，下一个结果用商环的结构刻画极大理想。

\begin{proposition}\label{pro:7.12}
假设 \(R\) 交换。理想 \(M\) 是极大理想当且仅当商环 \(R/M\) 是域。
\end{proposition}

\begin{proof}
这由格子同构定理连同命题 \ref{pro:7.9}(2) 推出。理想 \(M\) 极大当且仅当不存在满足 \(M \subsetneq I \subsetneq R\) 的理想 \(I\). 由格子同构定理，包含 \(M\) 的 \(R\) 的理想与 \(R/M\) 的理想一一对应，故 \(M\) 极大当且仅当 \(R/M\) 仅有的理想是 0 和 \(R/M\). 由命题 \ref{pro:7.9}(2) 我们看到 \(M\) 极大当且仅当 \(R/M\) 是域。
\end{proof}

上面的命题指出如何构造一些域：取任意带单位元的交换环 \(R\) 对 \(R\) 中极大理想的商。我们将在第 IV 部分用它，通过对 \(\mathbb{Z}[x]\) 取极大理想来构造所有有限域。

\begin{example}
（1）设 \(n\) 是非负整数。\(\mathbb{Z}\) 的理想 \(n\mathbb{Z}\) 是极大理想当且仅当 \(\mathbb{Z}/n\mathbb{Z}\) 是域。我们在 7.3 节看到这当且仅当 \(n\) 是素数。这也可以由上面例 2 中描述的 \(\mathbb{Z}\) 的理想的包含关系直接推出。

（2）理想 \((2, x)\) 是 \(\mathbb{Z}[x]\) 中的极大理想，因为它的商环是域 \(\mathbb{Z}/2\mathbb{Z}\)——参见上面例 3 以及 7.3 节末的例 5.

（3）\(\mathbb{Z}[x]\) 中的理想 \((x)\) 不是极大理想，因为 \((x) \subsetneq (2, x) \subsetneq \mathbb{Z}[x]\). 商环 \(\mathbb{Z}[x]/(x)\) 同构于 \(\mathbb{Z}\)（\(\mathbb{Z}[x]\) 中的理想 \((x)\) 是在 0 处求值给出的满环同态 \(\mathbb{Z}[x] \to \mathbb{Z}\) 的核）。由于 \(\mathbb{Z}\) 不是域，我们再次看到 \((x)\) 不是 \(\mathbb{Z}[x]\) 的极大理想。

（4）设 \(R\) 是所有从 \([0,1]\) 到 \(\mathbb{R}\) 的函数构成的环，对每个 \(a \in [0,1]\) 设 \(M_a\) 是在 \(a\) 处求值的核。由于求值是从 \(R\) 到 \(\mathbb{R}\) 的满同态，我们看到 \(R/M_a \cong \mathbb{R}\)，从而 \(M_a\) 是极大理想。类似地，在 \([0,1]\) 上的连续实值函数环中，在任意固定点处求值的核是极大理想。

（5）若 \(F\) 是域、\(G\) 是有限群，则增广理想 \(I\) 是群环 \(FG\) 的极大理想（参见前节末例 7）。增广理想是增广映射的核，后者是到域 \(F\) 上的满同态（即 \(FG/I \cong F\)，是域）。注意命题 \ref{pro:7.12} 不能直接应用，因为 \(FG\) 不必交换；然而命题 \ref{pro:7.12} 中“若 \(R/I\) 是域则 \(I\) 是极大理想”这一方向对任意环都成立。
\end{example}

\begin{definition}
假设 \(R\) 交换。理想 \(P\) 称为\textbf{素理想}，若 \(P \neq R\) 且只要两个元素 \(a, b \in R\) 之积 \(ab\) 是 \(P\) 的元素，就至少有一个 \(a, b\) 是 \(P\) 的元素。
\end{definition}

极大理想的概念相当直观，但素理想的定义可能显得有点奇怪。然而它是整数 \(\mathbb{Z}\) 中素数概念的自然推广。设 \(n\) 是非负整数。按上面的定义，理想 \(n\mathbb{Z}\) 是素理想，只要 \(n \neq 1\)（以保证理想真），且只要两个整数之积 \(ab\) 是 \(n\mathbb{Z}\) 的元素时至少有一个 \(a, b\) 是 \(n\mathbb{Z}\) 的元素。换句话说，若 \(n \neq 0\)，它必须具有性质：只要 \(n\) 整除 \(ab\)，\(n\) 必整除 \(a\) 或整除 \(b\). 这等价于通常的 \(n\) 是素数的定义。因此 \(\mathbb{Z}\) 的素理想正是由素数 \(p\) 生成的理想 \(p\mathbb{Z}\) 连同理想 0.

对整数 \(\mathbb{Z}\)，极大理想与非零素理想没有区别。一般并非如此，但我们不久将看到每个极大理想都是素理想。首先我们把素理想的概念翻译成商环的性质，如同我们在命题 \ref{pro:7.12} 中对极大理想所做的那样。回忆整环是带单位元 \(1 \neq 0\)、没有零因子的交换环。

\begin{proposition}\label{pro:7.13}
假设 \(R\) 交换。则理想 \(P\) 是 \(R\) 的素理想当且仅当商环 \(R/P\) 是整环。
\end{proposition}

\begin{proof}
这个证明只是把素理想的定义翻译成商的语言。理想 \(P\) 是素理想当且仅当 \(P \neq R\) 且只要 \(ab \in P\) 就有 \(a \in P\) 或 \(b \in P\). 用横线记号表示 \(R/P\) 的元素：\(\bar{r} = r+P\). 注意 \(r \in P\) 当且仅当元素 \(\bar{r}\) 在商环 \(R/P\) 中为零。因此用商的术语，\(P\) 是素理想当且仅当 \(R/P \neq \bar{0}\) 且只要 \(\overline{ab} = \bar{0}\) 就有 \(\bar{a} = \bar{0}\) 或 \(\bar{b} = \bar{0}\)，即 \(R/P\) 是整环。
\end{proof}

特别地推出带单位元的交换环是整环当且仅当 0 是素理想。

\begin{corollary}\label{cor:7.14}
假设 \(R\) 交换。\(R\) 的每个极大理想都是素理想。
\end{corollary}

\begin{proof}
若 \(M\) 是极大理想，则由命题 \ref{pro:7.12} \(R/M\) 是域。域是整环，故由命题 \ref{pro:7.13} 得推论。
\end{proof}

\begin{example}
（1）\(\mathbb{Z}\) 中由素数生成的主理想既是素理想又是极大理想。\(\mathbb{Z}\) 中的零理想是素理想但不是极大理想。

（2）理想 \((x)\) 是 \(\mathbb{Z}[x]\) 中的素理想，因为 \(\mathbb{Z}[x]/(x) \cong \mathbb{Z}\). 这个理想不是极大理想。理想 0 是 \(\mathbb{Z}[x]\) 中的素理想，但不是极大理想。
\end{example}

\subsection*{习题}
设 \(R\) 是带单位元 \(1 \neq 0\) 的环。

\begin{enumerate}
\item 设 \(L_j\) 是 \(M_n(R)\) 中第 \(j\) 列元素任意、其余元素为零的左理想，\(E_{ij}\) 是 \(M_n(R)\) 中 \((i, j)\) 位置为 1、其余为零的元素。证明对任意 \(i\) 有 \(L_j = M_n(R)E_{ij}\).〔见第 2 节习题 6。〕

\item 假设 \(R\) 交换。证明群环 \(RG\) 中的增广理想由 \(\{g-1 \mid g \in G\}\) 生成。证明若 \(G = \langle a \rangle\) 循环，则增广理想由 \(a-1\) 生成。

\item （a）设 \(p\) 是素数，\(G\) 是 \(p^n\) 阶阿贝尔群。证明群环 \(\mathbb{F}_pG\) 的幂零根是增广理想（参见第 3 节习题 29）。〔用前一习题。〕

（b）设 \(G = \{g_1, \ldots, g_n\}\) 是有限群，假设 \(R\) 交换。证明若 \(r\) 是 \(RG\) 中增广理想的任意元素，则 \(r(g_1 + \cdots + g_n) = 0\).〔用前一习题。〕

\item 假设 \(R\) 交换。证明 \(R\) 是域当且仅当 0 是极大理想。

\item 证明：若 \(M\) 是使 \(R/M\) 为域的理想，则 \(M\) 是极大理想（不要假设 \(R\) 交换）。

\item 证明 \(R\) 是除环当且仅当它仅有的左理想是 (0) 和 \(R\).（把“左”换成“右”类似结果成立。）

\item 设 \(R\) 是带 1 的交换环。证明多项式环 \(R[x]\) 中由 \(x\) 生成的主理想是素理想当且仅当 \(R\) 是整环。证明 \((x)\) 是极大理想当且仅当 \(R\) 是域。

\item 设 \(R\) 是整环。证明对元素 \(a, b \in R\)，\((a) = (b)\) 当且仅当对 \(R\) 的某个单位 \(u\) 有 \(a = ub\).

\item 设 \(R\) 是 \([0,1]\) 上所有连续函数构成的环，\(I\) 是 \(R\) 中满足 \(f(1/3) = f(1/2) = 0\) 的函数 \(f(x)\) 的集合。证明 \(I\) 是 \(R\) 的理想但不是素理想。

\item 假设 \(R\) 交换。证明：若 \(P\) 是 \(R\) 的素理想且 \(P\) 不含零因子，则 \(R\) 是整环。

\item 假设 \(R\) 交换。设 \(I, J\) 是 \(R\) 的理想，假设 \(P\) 是 \(R\) 的包含 \(IJ\) 的素理想（例如 \(P\) 包含 \(I \cap J\)）。证明 \(I\) 或 \(J\) 包含在 \(P\) 中。

\item 假设 \(R\) 交换，设 \(I = (a_1, a_2, \ldots, a_n)\)、\(J = (b_1, b_2, \ldots, b_m)\) 是 \(R\) 的两个有限生成理想。证明积理想 \(IJ\) 由元素 \(a_ib_j\)（\(i = 1, 2, \ldots, n\)，\(j = 1, 2, \ldots, m\)）有限生成。

\item 设 \(\varphi : R \to S\) 是交换环的同态。

（a）证明：若 \(P\) 是 \(S\) 的素理想，则或者 \(\varphi^{-1}(P) = R\)，或者 \(\varphi^{-1}(P)\) 是 \(R\) 的素理想。把它应用于 \(R\) 是 \(S\) 的子环、\(\varphi\) 是包含同态的特殊情形，推出若 \(P\) 是 \(S\) 的素理想则 \(P \cap R\) 或者等于 \(R\)，或者是 \(R\) 的素理想。

（b）证明：若 \(M\) 是 \(S\) 的极大理想且 \(\varphi\) 满射，则 \(\varphi^{-1}(M)\) 是 \(R\) 的极大理想。举一个 \(\varphi\) 不满射时这一结论不成立的例子。

\item 假设 \(R\) 交换。设 \(x\) 是未定元，\(f(x)\) 是 \(R[x]\) 中次数 \(n \geq 1\) 的首一多项式，用横线记号表示过渡到商环 \(R[x]/(f(x))\)。

（a）证明 \(R[x]/(f(x))\) 的每个元素都形如 \(\overline{p(x)}\)，其中 \(p(x)\) 是 \(R[x]\) 中次数小于 \(n\) 的多项式，即
\[
R[x]/(f(x)) = \left\{\overline{a_0 + a_1x + \cdots + a_{n-1}x^{n-1}} \ \middle|\ a_0, a_1, \ldots, a_{n-1} \in R\right\}.
\]
〔若 \(f(x) = x^n + b_{n-1}x^{n-1} + \cdots + b_0\)，则 \(\overline{x^n} = -(\overline{b_{n-1}x^{n-1} + \cdots + b_0})\)。用它约化商环中 \(x\) 的幂。〕

（b）证明：若 \(p(x)\) 和 \(q(x)\) 是 \(R[x]\) 中次数都小于 \(n\) 的不同多项式，则 \(\overline{p(x)} \neq \overline{q(x)}\).〔否则 \(p(x) - q(x)\) 是首一多项式 \(f(x)\) 的 \(R[x]\)-倍数。〕

（c）若 \(f(x) = a(x)b(x)\)，其中 \(a(x), b(x)\) 的次数都小于 \(n\)，证明 \(\overline{a(x)}\) 是 \(R[x]/(f(x))\) 中的零因子。

（d）若对某个幂零元素 \(a \in R\) 有 \(f(x) = x^n - a\)，证明 \(\overline{x}\) 在 \(R[x]/(f(x))\) 中幂零。

（e）设 \(p\) 是素数，假设 \(R = \mathbb{F}_p\) 且对某个 \(a \in \mathbb{F}_p\) 有 \(f(x) = x^p - a\). 证明 \(\overline{x-a}\) 在 \(R[x]/(f(x))\) 中幂零。〔用第 3 节习题 26(c)。〕

\item 设 \(x^2 + x + 1\) 是多项式环 \(\mathbb{F}_2[x]\) 的元素，用横线记号表示过渡到商环 \(E = \mathbb{F}_2[x]/(x^2+x+1)\).

（a）证明 \(E\) 有 4 个元素：\(\bar{0}, \bar{1}, \bar{x}\) 和 \(\overline{x+1}\).

（b）写出 \(E\) 的 \(4 \times 4\) 加法表，并由此推出加法群 \(E\) 同构于克莱因四元群。

（c）写出 \(E\) 的 \(4 \times 4\) 乘法表，证明 \(E \setminus \{\bar{0}\}\) 同构于 3 阶循环群。由此推出 \(E\) 是域。

\item 设 \(x^4 - 16\) 是多项式环 \(\mathbb{Z}[x]\) 的元素，用横线记号表示过渡到商环 \(E = \mathbb{Z}[x]/(x^4-16)\).

（a）求一个次数不超过 3、模 \((x^4-16)\) 同余于 \(7x^{13} - 11x^9 + 5x^5 - 2x^3 + 3\) 的多项式。

（b）证明 \(\overline{x-2}\) 和 \(\overline{x+2}\) 是 \(E\) 中的零因子。

\item 设 \(x^3 - 2x + 1\) 是多项式环 \(\mathbb{Z}[x]\) 的元素，用横线记号表示过渡到商环 \(E = \mathbb{Z}[x]/(x^3-2x+1)\). 设 \(p(x) = 2x^7 - 7x^5 + 4x^3 - 9x + 1\)，\(q(x) = (x-1)^4\).

（a）把 \(E\) 的下列每个元素表示成 \(\overline{f(x)}\) 的形式，其中 \(f(x)\) 是次数不超过 2 的多项式：\(\overline{p(x)}\)、\(\overline{q(x)}\)、\(\overline{p(x)+q(x)}\)、\(\overline{p(x)q(x)}\).

（b）证明 \(E\) 不是整环。

（c）证明 \(\overline{x}\) 是 \(E\) 中的单位。

\item 证明：若 \(R\) 是整环且 \(R[[x]]\) 是未定元 \(x\) 的形式幂级数环，则由 \(x\) 生成的主理想是素理想（参见第 2 节习题 3）。证明由 \(x\) 生成的主理想是极大理想当且仅当 \(R\) 是域。

\item 设 \(R\) 是带单位元的有限交换环。证明 \(R\) 的每个素理想都是极大理想。

\item 证明没有零因子的非零有限交换环是域（若环有单位元，这就是推论 \ref{cor:7.3}，故不要假设环有 1）。

\item 证明带单位元 \(1 \neq 0\)、没有零因子的有限环是域（你可以引用 Wedderburn 定理）。

\item 设 \(p \in \mathbb{Z}^+\) 是素数，定义 \(F_p\) \textbf{四元数}为
\[
a + bi + cj + dk, \quad a, b, c, d \in \mathbb{Z}/p\mathbb{Z},
\]
其中加法按分量，乘法使用与实四元数相同的 \(i, j, k\) 上的关系定义。

（a）证明 \(F_p\) 四元数是整四元数的同态像（参见第 1 节）。

（b）证明 \(F_p\) 四元数含零因子（因而不能是除环）。〔用前一习题。〕

\item 证明布尔环（参见第 1 节习题 15）中每个素理想都是极大理想。

\item 证明布尔环中每个有限生成理想都是主理想。

\item 假设 \(R\) 交换，且对每个 \(a \in R\) 存在整数 \(n > 1\)（依赖于 \(a\)）使 \(a^n = a\). 证明 \(R\) 的每个素理想都是极大理想。

\item 证明交换环 \(R\) 中的素理想包含每个幂零元素（参见第 1 节习题 13）。由此推出 \(R\) 的幂零根（参见第 3 节习题 29）包含在 \(R\) 的所有素理想的交中。（15.2 节将证明 \(R\) 的幂零根等于 \(R\) 的所有素理想的交。）

\item 设 \(R\) 是带单位元 \(1 \neq 0\) 的交换环。证明若 \(a\) 是 \(R\) 的幂零元素，则对所有 \(b \in R\)，\(1-ab\) 是单位。

\item 证明：若 \(R\) 是交换环且 \(N = (a_1, a_2, \ldots, a_m)\)，其中每个 \(a_i\) 是幂零元素，则 \(N\) 是幂零理想（参见第 3 节习题 37）。由此推出若 \(R\) 的幂零根有限生成，则它是幂零理想。

\item 设 \(p\) 是素数，\(G\) 是阶为 \(p\) 的幂的有限群（即 \(p\)-群）。证明群环 \(\mathbb{Z}/p\mathbb{Z}G\) 中的增广理想是幂零理想。（注意这个环可能不交换。）〔用习题 2。〕

\item 设 \(I\) 是交换环 \(R\) 的理想，定义
\[
\operatorname{rad} I = \{r \in R \mid r^n \in I \text{ 对某个 } n \in \mathbb{Z}^+\},
\]
称为 \(I\) 的\textbf{根}。证明 \(\operatorname{rad} I\) 是包含 \(I\) 的理想，且 \((\operatorname{rad} I)/I\) 是商环 \(R/I\) 的幂零根，即 \((\operatorname{rad} I)/I = \mathfrak{N}(R/I)\)（参见第 3 节习题 29）。

\item 交换环 \(R\) 的理想 \(I\) 称为\textbf{根理想}，若 \(\operatorname{rad} I = I\).

（a）证明 \(R\) 的每个素理想都是根理想。

（b）设 \(n > 1\) 是整数。证明 0 是 \(\mathbb{Z}/n\mathbb{Z}\) 中的根理想当且仅当 \(n\) 是互不相同素数的一次幂之积（即 \(n\) 无平方因子）。由此推出 \((n)\) 是 \(\mathbb{Z}\) 的根理想当且仅当 \(n\) 是 \(\mathbb{Z}\) 中互不相同素数的乘积。

\item 设 \(I\) 是交换环 \(R\) 的理想，定义
\[
\operatorname{Jac} I = \text{所有包含 } I \text{ 的 } R \text{ 的极大理想的交},
\]
约定 \(\operatorname{Jac} R = R\).（若 \(I\) 是零理想，\(\operatorname{Jac} 0\) 称为环 \(R\) 的\textbf{雅各布森根}，故 \(\operatorname{Jac} I\) 是 \(R/I\) 的雅各布森根在 \(R\) 中的原像。）

（a）证明 \(\operatorname{Jac} I\) 是 \(R\) 的包含 \(I\) 的理想。

（b）证明 \(\operatorname{rad} I \subseteq \operatorname{Jac} I\)，其中 \(\operatorname{rad} I\) 是习题 30 中定义的 \(I\) 的根。

（c）设 \(n > 1\) 是整数。用 \(n\) 的素因子分解描述 \(\operatorname{Jac} n\mathbb{Z}\).

\item 设 \(R\) 是所有从闭区间 \([0,1]\) 到 \(\mathbb{R}\) 的连续函数构成的环，对每个 \(c \in [0,1]\) 设 \(M_c = \{f \in R \mid f(c) = 0\}\)（回忆已证明 \(M_c\) 是 \(R\) 的极大理想）。

（a）证明若 \(M\) 是 \(R\) 的任意极大理想，则存在实数 \(c \in [0,1]\) 使 \(M = M_c\).

（b）证明若 \(b, c\) 是 \([0,1]\) 中的不同点，则 \(M_b \neq M_c\).

（c）证明 \(M_c\) 不等于由 \(x-c\) 生成的主理想。

（d）证明 \(M_c\) 不是有限生成理想。

\noindent 前一习题表明闭区间 \([0,1]\) 的点与 \([0,1]\) 上所有连续函数环 \(R\) 的极大理想的集合之间存在由 \(c \mapsto M_c\) 给出的双射。对 \(\mathbb{R}\) 的任意子集 \(X\)，或者更一般地对任意完全正则拓扑空间 \(X\)，映射 \(c \mapsto M_c\) 是从 \(X\) 到 \(R\) 的极大理想集合的单射，其中 \(R\) 是 \(X\) 上所有有界连续实值函数构成的环，\(M_c\) 是在 \(c\) 处为零的函数的极大理想。设 \(\beta(X)\) 是 \(R\) 的极大理想的集合。可以在 \(\beta(X)\) 上赋予一个拓扑，使得若把 \(X\) 与它在 \(\beta(X)\) 中的像等同，则 \(X\)（带其给定的拓扑）成为 \(\beta(X)\) 的子空间。此外，\(\beta(X)\) 在这个拓扑下是紧空间，称为 \(X\) 的 \textbf{Stone--Čech 紧化}。

\item 设 \(R\) 是所有从 \(\mathbb{R}\) 到 \(\mathbb{R}\) 的连续函数构成的环，对每个 \(c \in \mathbb{R}\) 设 \(M_c\) 是极大理想 \(\{f \in R \mid f(c) = 0\}\).

（a）设 \(I\) 是 \(R\) 中具有紧支集的函数 \(f(x)\) 的集合（即对充分大的 \(|x|\) 有 \(f(x) = 0\)）。证明 \(I\) 是 \(R\) 的理想但不是素理想。

（b）设 \(M\) 是 \(R\) 的（由（a）真地）包含 \(I\) 的极大理想。证明对任意 \(c \in \mathbb{R}\) 有 \(M \neq M_c\)（参见前一习题）。

\item 设 \(A = (a_1, a_2, \ldots, a_n)\) 是 \(R\) 的非零有限生成理想。证明存在理想 \(B\)，它在“不包含 \(A\)”这一性质下极大。〔用佐恩引理。〕

\item 假设 \(R\) 交换。证明 \(R\) 中素理想的集合（按包含关系）有极小元（可能是零理想）。〔用佐恩引理。〕

\item 交换环 \(R\) 称为\textbf{局部环}，若它有唯一的极大理想。证明：若 \(R\) 是极大理想为 \(M\) 的局部环，则 \(R - M\) 的每个元素都是单位。反之证明：若 \(R\) 是带 1 的交换环，其中非单位元集合构成理想 \(M\)，则 \(R\) 是极大理想为 \(M\) 的局部环。

\item 证明所有分母为奇数的有理数构成的环是局部环，其唯一的极大理想是由 2 生成的主理想。

\item 沿用第 1 节习题 26 的记号，设 \(K\) 是域，\(v\) 是 \(K\) 上的离散赋值，\(R\) 是 \(v\) 的赋值环。对每个整数 \(k \geq 0\) 定义
\[
A_k = \{r \in R \mid v(r) \geq k\} \cup \{0\}.
\]
（a）证明 \(A_k\) 是主理想且 \(A_0 \supseteq A_1 \supseteq A_2 \supseteq \cdots\).

（b）证明若 \(I\) 是 \(R\) 的任意非零理想，则对某个 \(k \geq 0\) 有 \(I = A_k\). 由此推出 \(R\) 是以 \(A_1\) 为唯一极大理想的局部环。

\item 假设 \(R\) 交换。证明下列命题等价：（另见第 1 节习题 13 和 14）

（i）\(R\) 恰有一个素理想；

（ii）\(R\) 的每个元素要么是幂零的要么是单位；

（iii）\(R/\mathfrak{N}(R)\) 是域（参见第 3 节习题 29）。

\item 交换环 \(R\) 的真理想 \(Q\) 称为\textbf{准素理想}，若只要 \(ab \in Q\) 且 \(a \notin Q\)，就存在正整数 \(n\) 使 \(b^n \in Q\).（注意这个定义中 \(a\) 与 \(b\) 的对称性蕴含：若 \(Q\) 是准素理想且 \(ab \in Q\) 而 \(a, b\) 都不在 \(Q\) 中，则 \(a\) 的某个正幂和 \(b\) 的某个正幂都在 \(Q\) 中。）建立关于准素理想的下列事实。

（a）\(\mathbb{Z}\) 的准素理想是 0 和 \((p^n)\)，其中 \(p\) 是素数、\(n\) 是正整数。

（b）\(R\) 的每个素理想都是准素理想。

（c）\(R\) 的理想 \(Q\) 是准素的当且仅当 \(R/Q\) 中每个零因子都是 \(R/Q\) 的幂零元素。

（d）若 \(Q\) 是准素理想，则 \(\operatorname{rad}(Q)\) 是素理想（参见习题 30）。
\end{enumerate}

\section{分式环}

本节自始至终 \(R\) 是交换环。命题 \ref{pro:7.2} 表明若 \(a\) 不为零也不是零因子，且在 \(R\) 中 \(ab = ac\)，则 \(b = c\). 因此不是零因子的非零元素享有单位的某些性质，而不必在 \(R\) 中有乘法逆元。另一方面，我们在 7.1 节看到零因子 \(a\) 不可能是 \(R\) 的单位，且按定义若 \(a\) 是零因子，我们不一定能在方程 \(ab = ac\) 中消去 \(a\) 得到 \(b = c\)（例如取 \(c = 0\)）。本节的目的是证明交换环 \(R\) 总是一个更大环 \(Q\) 的子环，其中 \(R\) 的每个不是零因子的非零元素都是 \(Q\) 中的单位。其主要应用是整环，此时这个环 \(Q\) 是域——称为它的\textbf{分式域}或\textbf{商域}。事实上，从 \(R\) 构造 \(Q\) 的范式正是从整环 \(\mathbb{Z}\) 构造有理数域 \(\mathbb{Q}\) 所提供的范式。

为看出从整环 \(\mathbb{Z}\) 构造 \(\mathbb{Q}\) 的本质特征，我们回顾分数的基本性质。每个有理数都可以用许多不同方式表示为两个整数之商（例如 \(-\frac12 = \frac{2}{-4} = \frac{-3}{6} = \cdots\)）。这些表示之间的关系是
\[
\frac{a}{b} = \frac{c}{d} \quad \text{当且仅当} \quad ad = bc.
\]
更确切地说，分数 \(\frac{a}{b}\) 是整数有序对 \((a, b)\)（\(b \neq 0\)）在等价关系“\((a, b) \sim (c, d)\) 当且仅当 \(ad = bc\)”下的等价类。分数的算术运算由
\[
\frac{a}{b} + \frac{c}{d} = \frac{ad+bc}{bd}, \qquad \frac{a}{b} \times \frac{c}{d} = \frac{ac}{bd}
\]
给出。这些是良定义的（与等价类代表元的选取无关），并使分数的集合成为交换环（事实上是域）\(\mathbb{Q}\). 整数 \(\mathbb{Z}\) 与 \(\mathbb{Q}\) 的子环 \(\{\frac{a}{1} \mid a \in \mathbb{Z}\}\) 等同，且每个非零整数 \(a\) 在 \(\mathbb{Q}\) 中有逆元 \(\frac{1}{a}\).

对任意交换环 \(R\) 尝试遵循相同的步骤似乎合理，允许任意分母。然而若 \(b\) 在 \(R\) 中为零或零因子，比如 \(bd = 0\)，且我们允许 \(b\) 作分母，则我们应当期待在“分式环”中有
\[
\frac{d}{1} = \frac{bd}{b} = \frac{0}{b} = 0
\]
（这里为方便起见假设 \(R\) 有 1）。因此若允许零或零因子作分母，必会发生某种坍缩，即不能期待 \(R\) 自然地作为这个“分式环”的子环出现。第二个限制更显然地由加法与乘法律施加：若环元素 \(b\) 和 \(d\) 允许作分母，则 \(bd\) 也必须作分母，即分母集合必须对 \(R\) 中乘法封闭。本节的主要结果表明这两个限制足以构造 \(R\) 的分式环。注意这个定理包含从 \(\mathbb{Z}\) 构造 \(\mathbb{Q}\) 作为特殊情形。

\begin{theorem}\label{thm:7.15}
设 \(R\) 是交换环。设 \(D\) 是 \(R\) 的任意非空子集，它不含 0、不含零因子且对乘法封闭（即对所有 \(a, b \in D\) 有 \(ab \in D\)）。则存在带 1 的交换环 \(Q\)，使 \(Q\) 以 \(R\) 为子环且 \(D\) 的每个元素都是 \(Q\) 中的单位。环 \(Q\) 具有下列额外性质。

（1）\(Q\) 的每个元素都形如 \(rd^{-1}\)，其中 \(r \in R\)、\(d \in D\). 特别地，若 \(D = R - \{0\}\)，则 \(Q\) 是域。

（2）（\(Q\) 的唯一性）环 \(Q\) 是包含 \(R\) 且其中 \(D\) 的所有元素都成为单位的“最小”环，意义如下。设 \(S\) 是任意带单位元的交换环，\(\varphi : R \to S\) 是任意单环同态，使得对每个 \(d \in D\)，\(\varphi(d)\) 是 \(S\) 中的单位。则存在单同态 \(\Phi : Q \to S\) 使 \(\Phi|_R = \varphi\). 换句话说，任何包含 \(R\) 的同构复制且其中 \(D\) 的所有元素都成为单位的环，也必包含 \(Q\) 的同构复制。
\end{theorem}

\textbf{注：}15.4 节将给出一个更一般的构造。一般结果的证明技术上更复杂，但依赖于与定理 \ref{thm:7.15} 的证明相同的基本原理和步骤。希望得到更大一般性的读者可以把下面的证明与 15.4 节的开头一起阅读。

\begin{proof}
设 \(\mathcal{F} = \{(r, d) \mid r \in R,\ d \in D\}\)，在 \(\mathcal{F}\) 上定义关系 \(\sim\) 为
\[
(r, d) \sim (s, e) \quad \text{当且仅当} \quad re = sd.
\]
立即看出这个关系是自反的、对称的。假设 \((r, d) \sim (s, e)\) 且 \((s, e) \sim (t, f)\). 则 \(re - sd = 0\) 且 \(sf - te = 0\). 把第一个等式乘以 \(f\)、第二个乘以 \(d\) 并相加得 \((rf - td)e = 0\). 由于 \(e \in D\) 既不为零也不是零因子，必有 \(rf - td = 0\)，即 \((r, d) \sim (t, f)\). 这证明 \(\sim\) 是传递的，从而是等价关系。把 \((r, d)\) 的等价类记作
\[
\frac{r}{d} = \{(a, b) \mid a \in R,\ b \in D,\ \text{且 } rb = ad\}.
\]
设 \(Q\) 是 \(\sim\) 下的等价类的集合。注意对所有 \(e \in D\) 有 \(\frac{r}{d} = \frac{re}{de}\)，因为 \(D\) 对乘法封闭。

我们现在在 \(Q\) 上定义加法与乘法结构：
\[
\frac{a}{b} + \frac{c}{d} = \frac{ad+bc}{bd}, \qquad \frac{a}{b} \times \frac{c}{d} = \frac{ac}{bd}.
\]
为证明 \(Q\) 是带单位元的交换环，有许多事情要验证：

（1）这些运算是良定义的（即不依赖于等价类代表元的选取）；

（2）\(Q\) 在加法下是阿贝尔群，其中加法单位元是对任意 \(d \in D\) 的 \(\frac{0}{d}\)，\(\frac{a}{b}\) 的加法逆元是 \(\frac{-a}{b}\)；

（3）乘法结合、分配且交换；

（4）\(Q\) 有单位元（对任意 \(d \in D\) 等于 \(\frac{d}{d}\)）。

这些都是完全直接的计算，只涉及 \(R\) 中的算术和 \(\sim\) 的定义。同样我们需要 \(D\) 对乘法封闭，才能使加法与乘法有定义。

例如为验证加法良定义，假设 \(\frac{a}{b} = \frac{a'}{b'}\)（即 \(ab' = a'b\)）且 \(\frac{c}{d} = \frac{c'}{d'}\)（即 \(cd' = c'd\)）。我们必须证明 \(\frac{ad+bc}{bd} = \frac{a'd'+b'c'}{b'd'}\)，即
\[
(ad+bc)(b'd') = (a'd'+b'c')(bd).
\]
这个等式左端是 \(ab'dd' + cb'bd'\)，把 \(ab'\) 替换为 \(a'b\)、把 \(cd'\) 替换为 \(c'd\) 得到 \(a'bdd' + c'dbb'\)，这正是右端。因此分数的加法良定义。验证（1）中其余部分和（2）至（4）的细节涉及更容易的操作，留作练习。

接下来我们通过定义
\[
\iota : R \to Q, \quad \iota : r \mapsto \frac{rd}{d}, \quad \text{其中 } d \text{ 是 } D \text{ 的任意元素}
\]
把 \(R\) 嵌入 \(Q\). 由于对所有 \(d, e \in D\) 有 \(\frac{rd}{d} = \frac{re}{e}\)，\(\iota(r)\) 不依赖于 \(d \in D\) 的选取。由于 \(D\) 对乘法封闭，可以直接验证 \(\iota\) 是环同态。此外 \(\iota\) 是单射，因为
\[
\iota(r) = \frac{0}{d} \iff \frac{rd}{d} = \frac{0}{d} \iff rd = 0 \iff r = 0,
\]
因为 \(d\)（因而 \(\frac{d}{d}\)）既不为零也不是零因子。\(Q\) 的子环 \(\iota(R)\) 因而同构于 \(R\). 我们今后把每个 \(r \in R\) 与 \(\iota(r)\) 等同，从而把 \(R\) 看作 \(Q\) 的子环。

接下来注意每个 \(d \in D\) 在 \(Q\) 中有乘法逆元：即若 \(d\) 由分数 \(\frac{de}{e}\) 表示，则它的乘法逆元是 \(\frac{e}{de}\). 于是每个 \(Q\) 的元素都可以写成 \(rd^{-1}\)，其中 \(r \in R\)、\(d \in D\). 特别地，若 \(D = R - \{0\}\)，则 \(Q\) 的每个非零元素都有乘法逆元，\(Q\) 是域。

剩下要建立 \(Q\) 的唯一性。假设 \(\varphi : R \to S\) 是单环同态，使得对所有 \(d \in D\)，\(\varphi(d)\) 是 \(S\) 中的单位。通过定义 \(\Phi(rd^{-1}) = \varphi(r)\varphi(d)^{-1}\)（对所有 \(r \in R\)、\(d \in D\)）把 \(\varphi\) 延拓为映射 \(\Phi : Q \to S\). 这个映射良定义，因为 \(rd^{-1} = se^{-1}\) 蕴含 \(re = sd\)，故 \(\varphi(r)\varphi(e) = \varphi(s)\varphi(d)\)，于是
\[
\Phi(rd^{-1}) = \varphi(r)\varphi(d)^{-1} = \varphi(s)\varphi(e)^{-1} = \Phi(se^{-1}).
\]
容易验证 \(\Phi\) 是环同态——细节留作练习。最后 \(\Phi\) 是单射，因为 \(rd^{-1} \in \ker\Phi\) 蕴含 \(r \in \ker\Phi \cap R = \ker\varphi\)；由于 \(\varphi\) 单射，这迫使 \(r\) 从而 \(rd^{-1}\) 为零。这完成了证明。
\end{proof}

\begin{definition}
设 \(R, D, Q\) 如定理 \ref{thm:7.15} 所述。

（1）环 \(Q\) 称为 \(D\) 关于 \(R\) 的\textbf{分式环}，记作 \(D^{-1}R\).

（2）若 \(R\) 是整环且 \(D = R - \{0\}\)，则 \(Q\) 称为 \(R\) 的\textbf{分式域}或\textbf{商域}。
\end{definition}

若 \(A\) 是域 \(F\) 的子集（例如 \(A\) 是 \(F\) 的子环），则 \(F\) 的所有包含 \(A\) 的子域的交是 \(F\) 的子域，称为 \(A\) \textbf{生成的子域}。这个子域是包含 \(A\) 的 \(F\) 的最小子域（即任何包含 \(A\) 的 \(F\) 的子域都包含 \(A\) 生成的子域）。

下一个推论表明包含整环 \(R\) 的最小的域是它的分式域。

\begin{corollary}\label{cor:7.16}
设 \(R\) 是整环，\(Q\) 是 \(R\) 的分式域。若域 \(F\) 含有一个同构于 \(R\) 的子环 \(R'\)，则 \(F\) 中由 \(R'\) 生成的子域同构于 \(Q\).
\end{corollary}

\begin{proof}
设 \(\varphi : R \to R' \subseteq F\) 是 \(R\) 到 \(R'\) 的（环）同构。特别地 \(\varphi : R \to F\) 是从 \(R\) 到域 \(F\) 的单同态。设 \(\Phi : Q \to F\) 是定理中 \(\varphi\) 到 \(Q\) 的延拓。由定理 \ref{thm:7.15}，\(\Phi\) 单射，故 \(\Phi(Q)\) 是 \(F\) 中 \(Q\) 的同构复制，包含 \(\varphi(R) = R'\). 现在，任何包含 \(R' = \varphi(R)\) 的 \(F\) 的子域都含元素 \(\varphi(r_1)\varphi(r_2)^{-1} = \varphi(r_1r_2^{-1})\)（对所有权 \(r_1, r_2 \in R\)）。由于 \(Q\) 的每个元素都形如 \(r_1r_2^{-1}\)（\(r_1, r_2 \in R\)），可知任何包含 \(R'\) 的 \(F\) 的子域都包含域 \(\Phi(Q)\)，故 \(\Phi(Q)\) 是 \(F\) 中由 \(R'\) 生成的子域，证得推论。
\end{proof}

\begin{example}
（1）若 \(R\) 是域，则其分式域就是 \(R\) 本身。

（2）整数 \(\mathbb{Z}\) 是分式域为有理数域 \(\mathbb{Q}\) 的整环。7.1 节的二次整数环 \(\mathcal{O}\) 是分式域为二次域 \(\mathbb{Q}(\sqrt{D})\) 的整环。

（3）\(\mathbb{Z}\) 的子环 \(2\mathbb{Z}\) 也没有零因子（但没有单位元）。它的分式域也是 \(\mathbb{Q}\). 注意单位元如何在分式域中“出现”。

（4）若 \(R\) 是任意整环，则多项式环 \(R[x]\) 也是整环。相应的分式域是 \(R\) 上变量 \(x\) 的\textbf{有理函数域}。这个域的元素形如 \(\frac{p(x)}{q(x)}\)，其中 \(p(x)\) 和 \(q(x)\) 是系数在 \(R\) 中的多项式，且 \(q(x)\) 不是零多项式。特别地，\(p(x)\) 和 \(q(x)\) 可以都是常多项式，故有理函数域含 \(R\) 的分式域：形如 \(\frac{a}{b}\) 的元素（\(a, b \in R\)，\(b \neq 0\)）。若 \(F\) 是域，我们用 \(F(x)\) 表示有理函数域。于是若 \(F\) 是整环 \(R\) 的分式域，则 \(R\) 上（因而也是 \(F\) 上）的有理函数域相同，即 \(F(x)\). 例如假设 \(R = \mathbb{Z}\)，故 \(F = \mathbb{Q}\). 若 \(p(x), q(x)\) 是 \(\mathbb{Q}[x]\) 中的多项式，则对某个整数 \(N\) 有 \(Np(x), Nq(x)\) 是整系数的（取 \(p(x), q(x)\) 中所有系数的一个公分母即可）。则 \(\frac{p(x)}{q(x)} = \frac{Np(x)}{Nq(x)}\) 可以写成两个整系数多项式之商，故 \(\mathbb{Q}[x]\) 的分式域与 \(\mathbb{Z}[x]\) 的分式域相同。

（5）若 \(R\) 是任意带单位元的交换环，\(d\) 在 \(R\) 中既不为零也不是零因子，我们可以通过取 \(D = \{1, d, d^2, d^3, \ldots\}\) 并定义 \(R[1/d]\) 为分式环 \(D^{-1}R\) 来构造环 \(R[1/d]\). 注意 \(R\) 是形如 \(\frac{r}{1}\) 的元素的子环。这样 \(R\) 中任何不是零因子的非零元素都可以在包含 \(R\) 的更大环中求逆。注意 \(R[1/d]\) 的元素形如系数在 \(R\) 中的 \(1/d\) 的多项式，这解释了记号。
\end{example}

\subsection*{习题}
设 \(R\) 是带单位元 \(1 \neq 0\) 的交换环。

\begin{enumerate}
\item 补全定理 \ref{thm:7.15} 证明中的所有细节。

\item 设 \(R\) 是整环，\(D\) 是 \(R\) 的闭于乘法的非空子集。证明分式环 \(D^{-1}R\) 同构于 \(R\) 的商域的某个子环（因而也是整环）。

\item 设 \(F\) 是域。证明 \(F\) 含有唯一的最小子域 \(F_0\)，且 \(F_0\) 同构于 \(\mathbb{Q}\) 或某个素数 \(p\) 的 \(\mathbb{Z}/p\mathbb{Z}\)（\(F_0\) 称为 \(F\) 的\textbf{素子域}）。〔见第 3 节习题 26。〕

\item 证明 \(\mathbb{R}\) 的任何子域都必含 \(\mathbb{Q}\).

\item 若 \(F\) 是域，证明 \(F[[x]]\)（系数在 \(F\) 中的未定元 \(x\) 的形式幂级数环）的分式域是形式洛朗级数环 \(F((x))\)（参见第 2 节习题 3 和 5）。证明幂级数环 \(\mathbb{Z}[[x]]\) 的分式域真包含在洛朗级数域 \(\mathbb{Q}((x))\) 中。〔考虑 \(e^x\) 的级数。〕

\item 证明实数 \(\mathbb{R}\) 含有一个子环 \(A\)，满足 \(1 \in A\)，且在“\(\frac{1}{2} \notin A\)”这一性质下（按包含关系）极大。〔用佐恩引理。〕（15.3 节习题 13 表明 \(\mathbb{R}\) 是 \(A\) 的商域，故 \(\mathbb{R}\) 是某个真子环的商域。）
\end{enumerate}

\section{中国剩余定理}

本节中除非另有说明，我们假设所有环都是带单位元 \(1 \neq 0\) 的交换环。

给定任意一族环（不必满足上述约定），它们的（环）\textbf{直积}定义为它们作为（阿贝尔）群的直积，并通过按分量定义乘法使其成为环。特别地，若 \(R_1\) 和 \(R_2\) 是两个环，我们用 \(R_1 \times R_2\) 表示它们的（作为环的）直积，即有序对 \((r_1, r_2)\)（\(r_1 \in R_1\)、\(r_2 \in R_2\)）的集合，其中加法与乘法按分量进行：
\[
(r_1, r_2) + (s_1, s_2) = (r_1+s_1,\ r_2+s_2), \quad (r_1, r_2)(s_1, s_2) = (r_1s_1,\ r_2s_2).
\]
注意从环 \(R\) 到直积环的映射 \(\varphi\) 是同态当且仅当到每个分量的诱导映射都是同态。

\(\mathbb{Z}\) 中两个整数 \(n\) 和 \(m\) 互素这一概念可以推广到任意环（甚至可以推广到未定义最大公因子概念的环）。在 \(\mathbb{Z}\) 中它等价于能在整数 \(x, y\) 中求解方程 \(nx + my = 1\)（这个事实在第 0 章陈述，将在第 8 章证明）。这又等价于作为理想 \(n\mathbb{Z} + m\mathbb{Z} = \mathbb{Z}\)（一般地 \(n\mathbb{Z} + m\mathbb{Z} = (m, n)\mathbb{Z}\)）。这引出下面的定义：

\begin{definition}
环 \(R\) 的理想 \(A\) 和 \(B\) 称为\textbf{互余的}，若 \(A + B = R\).
\end{definition}

回忆 \(R\) 的理想 \(A\) 和 \(B\) 的积 \(AB\) 是由形如 \(xy\)（\(x \in A\)、\(y \in B\)）的元素的全部有限和构成的理想（参见第 3 节习题 34）。若 \(A = (a)\)、\(B = (b)\)，则 \(AB = (ab)\). 更一般地，理想 \(A_1, A_2, \ldots, A_k\) 的积是形如 \(x_1x_2\cdots x_k\)（对所有 \(i\) 有 \(x_i \in A_i\)）的元素的全部有限和构成的理想。若 \(A_i = (a_i)\)，则 \(A_1\cdots A_k = (a_1\cdots a_k)\).

\begin{theorem}[中国剩余定理]\label{thm:7.17}
设 \(A_1, A_2, \ldots, A_k\) 是 \(R\) 的理想。映射
\[
R \to R/A_1 \times R/A_2 \times \cdots \times R/A_k, \quad \text{由 } r \mapsto (r+A_1,\ r+A_2,\ \ldots,\ r+A_k) \text{ 定义}
\]
是核为 \(A_1 \cap A_2 \cap \cdots \cap A_k\) 的环同态。若对每对 \(i, j \in \{1, 2, \ldots, k\}\)（\(i \neq j\)）理想 \(A_i\) 和 \(A_j\) 互余，则这个映射是满射且 \(A_1 \cap A_2 \cap \cdots \cap A_k = A_1A_2\cdots A_k\)，故
\[
R/(A_1A_2\cdots A_k) = R/(A_1 \cap A_2 \cap \cdots \cap A_k) \cong R/A_1 \times R/A_2 \times \cdots \times R/A_k.
\]
\end{theorem}

\begin{proof}
我们首先对 \(k = 2\) 证明；一般情形由归纳推出。设 \(A = A_1\)、\(B = A_2\). 考虑由
\[
\varphi : R \to R/A \times R/B, \quad \varphi(r) = (r \bmod A,\ r \bmod B)
\]
定义的映射，其中 \(r \bmod A\) 表示 \(R/A\) 中含 \(r\) 的类（即 \(r + A\)）。这个映射是环同态，因为 \(\varphi\) 对两个分量就是 \(R\) 到 \(R/A\) 和 \(R/B\) 的自然投影。\(\varphi\) 的核由所有既在 \(A\) 中又在 \(B\) 中的元素 \(r \in R\) 组成，即 \(A \cap B\). 为完成本情形的证明，还需证明当 \(A\) 和 \(B\) 互余时 \(\varphi\) 满射且 \(A \cap B = AB\).

由于 \(A + B = R\)，存在元素 \(x \in A\)、\(y \in B\) 使 \(x + y = 1\). 这个等式表明 \(\varphi(x) = (0, 1)\)、\(\varphi(y) = (1, 0)\)，因为例如 \(x\) 是 \(A\) 的元素且 \(x = 1 - y \in 1 + B\). 现在若 \((r_1 \bmod A, r_2 \bmod B)\) 是 \(R/A \times R/B\) 的任意元素，则元素 \(r_2x + r_1y\) 映到这个元素，因为
\[
\begin{aligned}
\varphi(r_2x + r_1y) &= \varphi(r_2)\varphi(x) + \varphi(r_1)\varphi(y) \\
&= (r_2 \bmod A,\ r_2 \bmod B)(0, 1) + (r_1 \bmod A,\ r_1 \bmod B)(1, 0) \\
&= (0,\ r_2 \bmod B) + (r_1 \bmod A,\ 0) = (r_1 \bmod A,\ r_2 \bmod B).
\end{aligned}
\]
这表明 \(\varphi\) 确实满射。最后，理想 \(AB\) 总是含于 \(A \cap B\). 若 \(A\) 和 \(B\) 互余且 \(x, y\) 如上，则对任意 \(c \in A \cap B\)，\(c = c\cdot 1 = cx + cy \in AB\). 这建立了反向包含 \(A \cap B \subseteq AB\)，并完成了 \(k = 2\) 时的证明。

一般情形容易由两个理想的情形归纳推出：取 \(A = A_1\)、\(B = A_2\cdots A_k\)，只要证明 \(A_1\) 与 \(A_2\cdots A_k\) 互余即可。由假设，对每个 \(i \in \{2, 3, \ldots, k\}\) 存在元素 \(x_i \in A_1\)、\(y_i \in A_i\) 使 \(x_i + y_i = 1\). 由于 \(x_i + y_i \equiv y_i \pmod{A_1}\)，可知 \(1 = (x_2+y_2)\cdots(x_k+y_k)\) 是 \(A_1 + (A_2\cdots A_k)\) 的元素。这完成了证明。
\end{proof}

这个定理的名字来自 \(m\) 与 \(n\) 互素时作为环的 \(\mathbb{Z}/mn\mathbb{Z} \cong (\mathbb{Z}/m\mathbb{Z}) \times (\mathbb{Z}/n\mathbb{Z})\) 这一特殊情形。我们先前只对加法群证明了这个同构。用数论的术语表述，这个同构与同时求解两个模互素整数的同余式有关（并断言这样的同余式总可以求解且唯一）。这类问题由古代中国人考虑，故得此名。习题中给出了一些例子。

由于中国剩余定理中的同构是环的同构，特别地两边的单位群必同构。容易看出环的直积中的单位是在每个坐标中都是单位的元素。对 \(\mathbb{Z}/mn\mathbb{Z}\) 的情形，中国剩余定理给出单位群的如下同构：
\[
(\mathbb{Z}/mn\mathbb{Z})^\times \cong (\mathbb{Z}/m\mathbb{Z})^\times \times (\mathbb{Z}/n\mathbb{Z})^\times.
\]
更一般地我们有以下结果。

\begin{corollary}\label{cor:7.18}
设 \(n\) 是正整数，\(n = p_1^{\alpha_1}p_2^{\alpha_2}\cdots p_k^{\alpha_k}\) 是其分解为互不相同素数的幂。则作为环
\[
\mathbb{Z}/n\mathbb{Z} \cong (\mathbb{Z}/p_1^{\alpha_1}\mathbb{Z}) \times (\mathbb{Z}/p_2^{\alpha_2}\mathbb{Z}) \times \cdots \times (\mathbb{Z}/p_k^{\alpha_k}\mathbb{Z}),
\]
特别地有乘法群的如下同构：
\[
(\mathbb{Z}/n\mathbb{Z})^\times \cong (\mathbb{Z}/p_1^{\alpha_1}\mathbb{Z})^\times \times (\mathbb{Z}/p_2^{\alpha_2}\mathbb{Z})^\times \times \cdots \times (\mathbb{Z}/p_k^{\alpha_k}\mathbb{Z})^\times.
\]
\end{corollary}

若比较最后这个同构两边的阶，我们得到欧拉 \(\varphi\)-函数的公式
\[
\varphi(n) = \varphi(p_1^{\alpha_1})\varphi(p_2^{\alpha_2})\cdots\varphi(p_k^{\alpha_k}).
\]
这又蕴含 \(\varphi\) 是初等数论中所称的\textbf{乘性函数}，即只要 \(a\) 与 \(b\) 是互素的正整数就有 \(\varphi(ab) = \varphi(a)\varphi(b)\). \(\varphi\) 在素数幂 \(p^\alpha\) 上的值容易看出是 \(\varphi(p^\alpha) = p^{\alpha-1}(p-1)\)（参见第 0 章）。由这一点以及 \(\varphi\) 的乘性，我们得到它在所有正整数上的值。

推论 \ref{cor:7.18} 也是确定阿贝尔群 \((\mathbb{Z}/n\mathbb{Z})^\times\) 分解为循环群直积的一步。完整的结构在 9.5 节末导出。

\subsection*{习题}
设 \(R\) 是带单位元 \(1 \neq 0\) 的环。

\begin{enumerate}
\item 元素 \(e \in R\) 称为\textbf{幂等元}，若 \(e^2 = e\). 假设 \(e\) 是 \(R\) 中的幂等元且对所有 \(r \in R\) 有 \(er = re\). 证明 \(Re\) 和 \(R(1-e)\) 是 \(R\) 的双边理想且 \(R \cong Re \times R(1-e)\). 证明 \(e\) 和 \(1-e\) 分别是子环 \(Re\) 和 \(R(1-e)\) 的单位元。

\item 设 \(R\) 是带单位元 \(1 \neq 0\) 的有限布尔环（参见第 1 节习题 15）。证明 \(R \cong \mathbb{Z}/2\mathbb{Z} \times \mathbb{Z}/2\mathbb{Z} \times \cdots \times \mathbb{Z}/2\mathbb{Z}\).〔用前一习题。〕

\item 设 \(R\) 和 \(S\) 是带单位元的环。证明 \(R \times S\) 的每个理想都形如 \(I \times J\)，其中 \(I\) 是 \(R\) 的理想、\(J\) 是 \(S\) 的理想。

\item 证明：若 \(R\) 和 \(S\) 是非零环，则 \(R \times S\) 永远不是域。

\item 设 \(n_1, n_2, \ldots, n_k\) 是两两互素的整数：对所有 \(i \neq j\) 有 \((n_i, n_j) = 1\).

（a）证明中国剩余定理蕴含：对任意 \(a_1, \ldots, a_k \in \mathbb{Z}\)，同余方程组
\[
x \equiv a_1 \pmod{n_1},\quad \ldots,\quad x \equiv a_k \pmod{n_k}
\]
在 \(\mathbb{Z}\) 中有解 \(x\)，且解 \(x\) 模 \(n = n_1n_2\cdots n_k\) 唯一。

（b）设 \(n_i' = n/n_i\) 是 \(n\) 除以 \(n_i\) 之商，由假设它与 \(n_i\) 互素。设 \(t_i\) 是 \(n_i'\) 模 \(n_i\) 的逆元。证明（a）中的解 \(x\) 由
\[
x \equiv a_1t_1n_1' + a_2t_2n_2' + \cdots + a_kt_kn_k' \pmod{n}
\]
给出。注意元素 \(t_i\) 可以用预备知识第 2 节描述的欧几里得算法快速求出（写出 \(an_i + bn_i' = (n_i, n_i') = 1\) 即得 \(t_i = b\)），而这些随后快速给出上述同余方程组对任意 \(a_1, a_2, \ldots, a_k\) 的解。

（c）求解同余方程组
\[
x \equiv 1 \pmod 8,\quad x \equiv 2 \pmod{25},\quad x \equiv 3 \pmod{81}
\]
以及同余方程组
\[
y \equiv 5 \pmod 8,\quad y \equiv 12 \pmod{25},\quad y \equiv 47 \pmod{81}.
\]

\item 设 \(f_1(x), f_2(x), \ldots, f_k(x)\) 是次数相同（都为 \(d\)）的整系数多项式。设 \(n_1, n_2, \ldots, n_k\) 是两两互素的整数（即对所有 \(i \neq j\) 有 \((n_i, n_j) = 1\)）。用中国剩余定理证明存在次数为 \(d\) 的整系数多项式 \(f(x)\) 使
\[
f(x) \equiv f_1(x) \pmod{n_1},\quad f(x) \equiv f_2(x) \pmod{n_2},\quad \ldots,\quad f(x) \equiv f_k(x) \pmod{n_k},
\]
即 \(f(x)\) 的系数与 \(f_i(x)\) 的系数模 \(n_i\) 相同。证明若所有 \(f_i(x)\) 都是首一的，则 \(f(x)\) 也可以取为首一的。〔对每个系数分别应用 \(\mathbb{Z}\) 中的中国剩余定理。〕

\item 设 \(m, n\) 是正整数且 \(n\) 整除 \(m\)。证明自然满环投影 \(\mathbb{Z}/m\mathbb{Z} \to \mathbb{Z}/n\mathbb{Z}\) 在单位上也是满射：\((\mathbb{Z}/m\mathbb{Z})^\times \to (\mathbb{Z}/n\mathbb{Z})^\times\).

\noindent 下面四个习题发展正极限的概念以及“对偶”的反极限概念。这些习题中 \(I\) 是带偏序 \(\leq\) 的非空指标集（参见附录 I）。对每个 \(i \in I\)，\(A_i\) 是加法阿贝尔群。习题 8 中还假设 \(I\) 是\textbf{有向集}：对任意 \(i, j \in I\) 存在 \(k \in I\) 使 \(i \leq k\) 且 \(j \leq k\).

\item 假设对每对满足 \(i \leq j\) 的指标 \(i, j\) 有映射 \(p_{ij} : A_i \to A_j\)，满足：

i. 只要 \(i \leq j \leq k\) 就有 \(p_{jk} \circ p_{ij} = p_{ik}\)，且

ii. 对所有 \(i \in I\) 有 \(p_{ii} = 1\).

设 \(B\) 是所有 \(A_i\) 的无交并。在 \(B\) 上定义关系 \(\sim\) 为：对 \(a \in A_i\)、\(b \in A_j\)，
\[
a \sim b \quad \text{当且仅当存在满足 } i, j \leq k \text{ 的 } k \text{ 使 } p_{ik}(a) = p_{jk}(b).
\]
（a）证明 \(\sim\) 是 \(B\) 上的等价关系。（等价类的集合称为有向系 \(\{A_i\}\) 的\textbf{正极限}或\textbf{归纳极限}，记作 \(\varinjlim A_i\). 在本习题其余部分设 \(A = \varinjlim A_i\).）

（b）用 \(\bar{x}\) 表示 \(x\) 在 \(A\) 中的类，定义 \(p_i : A_i \to A\) 为 \(p_i(a) = \bar{a}\). 证明若每个 \(p_{ij}\) 都是单射，则所有 \(p_i\) 都是单射（于是我们可以把每个 \(A_i\) 与 \(A\) 的子集等同）。

（c）假设所有 \(p_{ij}\) 都是群同态。对 \(a \in A_i\)、\(b \in A_j\) 证明运算
\[
\bar{a} + \bar{b} = \overline{p_{ik}(a) + p_{jk}(b)}
\]
（其中 \(k\) 是任意满足 \(i, j \leq k\) 的指标）良定义并使 \(A\) 成为阿贝尔群。由此推出（b）中的映射 \(p_i\) 是从 \(A_i\) 到 \(A\) 的群同态。

（d）证明：若所有 \(A_i\) 都是带 1 的交换环且所有 \(p_{ij}\) 都是把 1 送到 1 的环同态，则 \(A\) 同样可以被赋予带 1 的交换环结构，使所有 \(p_i\) 都是环同态。

（e）在（c）的假设下证明正极限具有如下泛性质：若 \(C\) 是任意阿贝尔群，使得对每个 \(i \in I\) 存在同态 \(\varphi_i : A_i \to C\)，且只要 \(i \leq j\) 就有 \(\varphi_i = \varphi_j \circ p_{ij}\)，则存在唯一的同态 \(\varphi : A \to C\) 使对所有 \(i\) 有 \(\varphi \circ p_i = \varphi_i\).

\item 设 \(I\) 是实轴上包含固定实数 \(p\) 的所有开区间 \(U = (a, b)\) 的集合。按反向包含对它们排序：\(U \geq V\) 若 \(V \subseteq U\)（注意此时 \(I\) 是有向集）。对每个 \(U\) 设 \(A_U\) 是 \(U\) 上连续实值函数构成的环。对 \(V \subseteq U\) 定义限制映射 \(\rho_{UV} : A_U \to A_V\) 为 \(f \mapsto f|_V\)，即函数在 \(U\) 上的通常限制到子集 \(V\) 上（容易看出这是环同态）。设 \(A = \varinjlim A_U\) 是正极限。用前一习题的记号，证明映射 \(p_U : A_U \to A\) 都不是单射但都是满射（\(A\) 称为在 \(p\) 处连续函数的\textbf{芽环}）。

\noindent 我们现在发展反极限的概念。继续假设 \(I\) 是偏序集（但不必有向），对所有 \(i \in I\) 有群 \(A_i\).

\item 假设对每对满足 \(i \geq j\) 的指标 \(i, j\) 有映射 \(\pi_{ji} : A_j \to A_i\)，满足：

i. 只要 \(i \geq j \geq k\) 就有 \(\pi_{ji} \circ \pi_{kj} = \pi_{ki}\)，且

ii. 对所有 \(i \in I\) 有 \(\pi_{ii} = 1\).

设 \(P\) 是直积 \(\prod_{i \in I} A_i\) 中满足只要 \(i \geq j\) 就有 \(\pi_{ji}(a_j) = a_i\) 的元素 \((a_i)_{i \in I}\) 的集合（这里 \(a_i\) 和 \(a_j\) 分别是直积中元素的第 \(i\) 和第 \(j\) 个分量）。集合 \(P\) 称为系 \(\{A_i\}\) 的\textbf{反极限}或\textbf{射影极限}，记作 \(\varprojlim A_i\).

（a）假设所有 \(\pi_{ji}\) 都是群同态。证明 \(P\) 是直积群的子群（参见第 5.1 节习题 15）。

（b）假设（a）中的假设成立，令 \(I = \mathbb{Z}^+\)（通常序）。对每个 \(i \in I\) 设 \(\Pi_i : P \to A_i\) 是 \(P\) 到其第 \(i\) 个分量的投影。证明若每个 \(\pi_{ji}\) 都是满射，则所有 \(\Pi_i\) 也是满射（于是每个 \(A_i\) 都是 \(P\) 的商群）。

（c）证明：若所有 \(A_i\) 都是带 1 的交换环且所有 \(\pi_{ji}\) 都是把 1 送到 1 的环同态，则 \(P\) 同样可以被赋予带 1 的交换环结构，使所有 \(\Pi_i\) 都是环同态。

（d）在（a）的假设下证明反极限具有如下泛性质：若 \(D\) 是任意群，使得对每个 \(i \in I\) 存在同态 \(\pi_i : D \to A_i\) 且只要 \(i \geq j\) 就有 \(\pi_i = \pi_{ji}\circ\pi_j\)，则存在唯一的同态 \(\pi : D \to P\) 使对所有 \(i\) 有 \(\Pi_i \circ \pi = \pi_i\).

\item 设 \(p\) 是素数，令 \(I = \mathbb{Z}^+\)，\(A_i = \mathbb{Z}/p^i\mathbb{Z}\)，\(\pi_{ji}\) 是自然投影映射 \(\pi_{ji} : a \pmod{p^j} \mapsto a \pmod{p^i}\). 反极限 \(\varprojlim \mathbb{Z}/p^i\mathbb{Z}\) 称为 \textbf{\(p\)-进整数环}，记作 \(\mathbb{Z}_p\).

（a）证明 \(\mathbb{Z}_p\) 的每个元素可以唯一地写成无穷形式和 \(b_0 + b_1p + b_2p^2 + b_3p^3 + \cdots\)，其中每个 \(b_i \in \{0, 1, \ldots, p-1\}\). 描述对应于环 \(\mathbb{Z}_p\) 中加法与乘法的这类形式和的加法与乘法规则。〔把每个 \(\mathbb{Z}/p^i\mathbb{Z}\) 中的最小剩余写成 \(p\) 进制展开，然后描述映射 \(\pi_{ji}\).〕（特别要注意 \(\mathbb{Z}_p\) 是不可数的。）

（b）证明 \(\mathbb{Z}_p\) 是整环且含有整数的一个复制。

（c）证明（a）中的 \(b_0 + b_1p + b_2p^2 + b_3p^3 + \cdots\) 是 \(\mathbb{Z}_p\) 中的单位当且仅当 \(b_0 \neq 0\).

（d）证明 \(p\mathbb{Z}_p\) 是 \(\mathbb{Z}_p\) 的唯一极大理想且 \(\mathbb{Z}_p/p\mathbb{Z}_p \cong \mathbb{Z}/p\mathbb{Z}\)（其中 \(p = 0 + 1\cdot p + 0\cdot p^2 + 0\cdot p^3 + \cdots\)）。证明 \(\mathbb{Z}_p\) 的每个理想都形如 \(p^n\mathbb{Z}_p\)（\(n \geq 0\) 为某个整数）。

（e）证明若 \(a_1 \not\equiv 0 \pmod p\)，则存在 \(\mathbb{Z}_p\) 中的元素 \(\alpha = (\alpha_i)\) 满足 \(\alpha_i^{p-1} = 1 \pmod{p^i}\) 且对所有 \(j\) 有 \(\pi_{ji}(\alpha_j) = \alpha_i\). 由此推出 \(\mathbb{Z}_p\) 含有 \(p-1\) 个互不相同的 \((p-1)\) 次单位根。
\end{enumerate}
